\chapter{Memory Slots and Lookup Memory}\label{ch:lookup}

The memories of the previous chapters are \emph{computed}. A KV cache is a stack of activations that
must be recomputed for every new context, a recurrent state is a compressed summary that is
overwritten token by token, and the fast weights of \cref{ch:ttt} are written by gradient steps.
None of them has a notion of \emph{where} a piece of knowledge lives. This chapter studies the
opposite design: memories made of individually \emph{addressable} entries --- slots, table rows,
stored token sequences --- that persist beyond a single forward pass and are reached through an
explicit address computation. The survey of \citet{zhoubian2026memory} groups such designs under
explicit memory, in its sections on dedicated memory slots (\S IV-A(d)), lookup-based and
retrieval-oriented memory (\S IV-B) and conditional parameter memory (\S IV-C).

Addressable memory matters for three reasons. First, capacity can be decoupled from compute: a
table with a billion rows costs the same per token as a table with a thousand rows if only a
constant number of rows is read. Second, an entry that has an address can be inspected, edited or
moved to another level of the hardware memory hierarchy. Third, it gives the network a native
primitive for \emph{knowledge lookup}, which a standard Transformer can only imitate by spending
layers of attention and feed-forward computation on it \citep{cheng2026engram}.

\Cref{sec:c10-background} builds the background from scratch: the read/write interface of an
addressable memory, nearest-neighbour datastores (kNN-LM), the view of feed-forward layers as
key--value memories, product-key memory with its $\bigO(\sqrt{N})$ lookup, hashed $N$-gram
embeddings, Mixture-of-Experts as conditional parameter memory, and two 2023--2024 precursors.
The four systems of the chapter then follow, ordered from the most ``recurrent'' to the most
``static'': LM2, a gated bank of memory slots inside every decoder layer
(\cref{sec:c10-lm2}); Engram, DeepSeek's hashed $N$-gram tables that act as a new axis of
sparsity next to MoE (\cref{sec:c10-engram}); ExplicitLM, a bank of human-readable token
sequences reached by a two-stage differentiable retrieval (\cref{sec:c10-explicitlm}); and
MemoryLLM, which turns the feed-forward layers themselves into context-free per-token lookup
tables (\cref{sec:c10-memoryllm}). \Cref{sec:c10-comparison} compares them.

% =====================================================================================
\section{Background: addressable memory inside a language model}\label{sec:c10-background}

\subsection{A common read/write interface}\label{sec:c10-interface}

All memories in this chapter fit one template. A memory holds $N$ \emph{entries}. Entry $i$ has an
address (a key $\vk_i$, an integer index, or both) and a content $\vv_i \in \R^{d_v}$. Reading
involves three steps:
\begin{enumerate}
  \item \textbf{Address.} From the current context (token ids, a hidden state $\vh_t$, or both)
        compute a small set of indices $\mathcal{I}_t \subset \{1,\dots,N\}$ with
        $\abs{\mathcal{I}_t} = k \ll N$, and weights $w_{t,i}$ for $i \in \mathcal{I}_t$.
  \item \textbf{Read.} Combine the selected contents, usually
        $\vr_t = \sum_{i \in \mathcal{I}_t} w_{t,i}\, \vv_i$.
  \item \textbf{Fuse.} Inject $\vr_t$ into the network, typically by adding a gated or projected
        version of it to the residual stream.
\end{enumerate}
Writing changes entries: by gradient descent during training (\emph{offline} memory, in the
survey's terminology), or by an explicit rule at inference time (\emph{online} memory).

\begin{definition}[Explicit, addressable memory]\label{def:c10-explicit}
Following \citet{zhoubian2026memory} (\S IV-B(g)), a memory is \emph{explicit} when (i) its entries
are individually addressable rather than compressed into one evolving state, (ii) its capacity
scales with the number of entries rather than with a hidden dimension, and (iii) its entries
persist independently of the forward computation that reads them.
\end{definition}

The KV cache satisfies (i) but not (iii): its entries are activations of the current context and
disappear with it. A recurrent state violates (i) and (ii). The memories below satisfy all three,
with one interesting exception: the memory bank of LM2 is a per-sequence state that is
nevertheless organized into slots.

The cost of a memory is dominated by the address step, and the address step is where the designs
differ. \Cref{tab:c10-addressing} lists the schemes used in this chapter; the rest of this section
derives each of them.

\begin{table}[t]
\centering
\small
\caption{Addressing schemes for a memory of $N$ entries with key dimension $d_k$, reading $k$
entries. ``Static'' means the address depends only on token ids and can be computed before the
forward pass.}
\label{tab:c10-addressing}
\begin{tabularx}{\linewidth}{@{}lXll@{}}
\toprule
Scheme & Address computation & Address cost & Static? \\
\midrule
Dense attention over slots & $\softmax(\vq\T\mK\T/\sqrt{d_k})$ over all $N$ keys & $\bigO(N d_k)$ & no \\
Exact nearest neighbours & $k$ keys closest to $\vq$ & $\bigO(N d_k)$ & no \\
Approximate nearest neighbours & index structure (clusters, graphs) & sublinear in $N$ & no \\
Product keys & two half-queries against $\sqrt{N}$ sub-keys each & $\bigO(\sqrt{N} d_k + k^2)$ & no \\
MoE router & top-$k$ of $E$ router logits & $\bigO(E d)$ & no \\
Hashing of token ids & arithmetic on the last $n$ token ids & $\bigO(n)$ & yes \\
Token id (embedding table) & the id itself & $\bigO(1)$ & yes \\
\bottomrule
\end{tabularx}
\end{table}

\subsection{Datastore memory: the kNN-LM}\label{sec:c10-knnlm}

The cleanest example of explicit memory attached to a language model is the $k$-nearest-neighbour
language model of \citet{khandelwal2020knnlm}. Take a trained Transformer LM and let
$f(c) \in \R^{d}$ be its final hidden state after reading a context $c$ (the vector that is fed to
the output softmax). Run the model once over a corpus and store one pair per token position,
\begin{equation}
  \mathcal{D} = \bigl\{ \bigl(f(c_i),\, w_i\bigr) \bigr\}_{i=1}^{N},
  \label{eq:c10-datastore}
\end{equation}
where $c_i$ is the prefix before position $i$ and $w_i$ the token that followed it. The keys are
neural representations; the values are plain tokens. At test time, with context $x$, the model
computes $\vq = f(x)$, retrieves the set $\mathcal{N}$ of the $k$ pairs whose keys are nearest to
$\vq$ under a distance $\delta$ (squared Euclidean distance in the original work), and forms a
distribution over the vocabulary by a softmax over negative distances:
\begin{equation}
  p_{\mathrm{kNN}}(y \mid x) = \frac{\sum_{(\vk_i, w_i) \in \mathcal{N}} \mathbb{1}[y = w_i]\,
  \exp\bigl(-\delta(\vk_i, \vq)\bigr)}{\sum_{(\vk_i, w_i) \in \mathcal{N}}
  \exp\bigl(-\delta(\vk_i, \vq)\bigr)} .
  \label{eq:c10-knn}
\end{equation}
The final prediction interpolates the two models with a scalar $\lambda \in [0,1]$:
$p(y \mid x) = \lambda\, p_{\mathrm{kNN}}(y \mid x) + (1-\lambda)\, p_{\mathrm{LM}}(y \mid x)$.

\begin{example}[Three neighbours]
Suppose the three nearest keys have distances $1, 2, 4$ and values \emph{Paris}, \emph{Paris},
\emph{Lyon}. Then $p_{\mathrm{kNN}}(\textit{Paris}) = (e^{-1}+e^{-2})/(e^{-1}+e^{-2}+e^{-4})
= (0.3679 + 0.1353)/0.5215 \approx 0.965$ and $p_{\mathrm{kNN}}(\textit{Lyon}) \approx 0.035$.
Every other token gets zero probability from the datastore, which is why the interpolation with
$p_{\mathrm{LM}}$ is needed.
\end{example}

The survey treats kNN-LM as the canonical explicit datastore (\S IV-B(a)): the datastore is not
part of the forward computation graph, survives across inference calls, and can be enlarged or
edited without touching the network weights. Its weaknesses point at the designs of this chapter.
The memory is \emph{external} to training (the network never learns to use it), and the address
step is a nearest-neighbour search over $N$ keys of dimension $d$. Exact search costs
$\bigO(Nd)$ per token; practical systems use approximate indices. Memorizing Transformers
\citep{wu2022memorizing} and LongMem \citep{wang2023longmem} move similar retrieval inside the
network, but store activations of past inputs; the survey treats them as a boundary class between
explicit datastores and implicit activation memory (\S IV-B(b)).

\subsection{Feed-forward layers as key--value memories}\label{sec:c10-ffnkv}

A Transformer block (\cref{ch:foundations}) contains an attention sub-layer and a position-wise
feed-forward network (FFN). Drop the biases and write the FFN with weights
$\mW_1 \in \R^{d_{\mathrm{ff}} \times d}$ and $\mW_2 \in \R^{d \times d_{\mathrm{ff}}}$ and an
elementwise nonlinearity $\phi$:
\begin{equation}
  \mathrm{FFN}(\vx) = \mW_2\, \phi(\mW_1 \vx).
\end{equation}
Call the $i$-th row of $\mW_1$ a key $\vk_i\T$ and the $i$-th column of $\mW_2$ a value $\vv_i$.
Expanding the matrix product column by column gives
\begin{equation}
  \mathrm{FFN}(\vx) = \sum_{i=1}^{d_{\mathrm{ff}}} \phi\bigl(\vk_i\T \vx\bigr)\, \vv_i .
  \label{eq:c10-ffn-kv}
\end{equation}
This has exactly the read form $\vr = \sum_i w_i \vv_i$ of \cref{sec:c10-interface}, with
$N = d_{\mathrm{ff}}$ slots and unnormalized weights $w_i = \phi(\vk_i\T\vx)$ instead of softmax
weights. \citet{c01-geva2021ffn} made the interpretation empirical: keys respond to recognizable
patterns in the input prefix, and values promote particular output tokens when projected onto the
vocabulary. The gated variant used in modern models,
$\mW_2(\SiLU(\mW_1\vx) \had \mW_3\vx)$, has the same structure with data-dependent value scaling.

Equation~\eqref{eq:c10-ffn-kv} explains why FFNs are often called the knowledge store of a
Transformer, and also why that store is \emph{implicit}: (i) the address $\vx$ is the full
contextualized hidden state, so it is unclear which ``slot'' a given input token will touch;
(ii) every slot is evaluated for every token, so the read costs $\bigO(d_{\mathrm{ff}} d)$ and
capacity can only grow with compute; (iii) knowledge is spread over many slots and layers. There
are two ways out, and this chapter covers both: make the address \emph{sparse}, so that only a few
of a very large number of slots are touched (product keys, MoE, Engram, ExplicitLM), or make the
address \emph{context-free}, so that the slot read by a token is known in advance (Engram's hashed
token ids, MemoryLLM).

\subsection{Product-key memory and its \texorpdfstring{$\bigO(\sqrt{N})$}{O(sqrt N)} lookup}\label{sec:c10-pkm}

Suppose we want a key--value layer with $N$ slots, $N$ in the millions, read by top-$k$ attention:
score every key against a query $\vq \in \R^{d_k}$, keep the $k$ best, and average their values.
Scoring all keys costs $\bigO(N d_k)$, which defeats the purpose. The product-key memory (PKM) of
\citet{c10-lample2019pkm} removes this cost by giving the keys a product structure.

\paragraph{Construction.} Let $N = n^2$. Keep two small sets of \emph{sub-keys},
$\mC^{(1)} = [\bm{c}^{(1)}_1, \dots, \bm{c}^{(1)}_n]\T$ and $\mC^{(2)} = [\bm{c}^{(2)}_1, \dots, \bm{c}^{(2)}_n]\T$,
each $\bm{c} \in \R^{d_k/2}$. Slot $(i,j)$, for $i,j \in \{1,\dots,n\}$, has the concatenated key
$\vk_{ij} = [\bm{c}^{(1)}_i ; \bm{c}^{(2)}_j] \in \R^{d_k}$. The $N$ keys are never stored: only the
$2n = 2\sqrt{N}$ sub-keys are. Split the query the same way, $\vq = [\vq^{(1)} ; \vq^{(2)}]$. Because
the inner product of concatenations is the sum of the inner products of the halves,
\begin{equation}
  s_{ij} \defeq \vq\T \vk_{ij} = \underbrace{\vq^{(1)\top} \bm{c}^{(1)}_i}_{a_i}
  + \underbrace{\vq^{(2)\top} \bm{c}^{(2)}_j}_{b_j} .
  \label{eq:c10-pkm-score}
\end{equation}
So the $N$ scores form an $n \times n$ outer sum $s_{ij} = a_i + b_j$ of two vectors that cost
only $2n$ half-length dot products to compute.

\paragraph{Exact top-$k$ from two short lists.} Let $\mathcal{I}_1 = \TopK(\va, k)$ and
$\mathcal{I}_2 = \TopK(\vb, k)$ be the indices of the $k$ largest entries of $\va$ and $\vb$.

\begin{proposition}[Product-key search is exact]\label{prop:c10-pkm}
Assume the scores $s_{ij}$ are pairwise distinct. Every pair $(i,j)$ among the $k$ largest
scores of the full $n \times n$ grid satisfies $i \in \mathcal{I}_1$ and $j \in \mathcal{I}_2$.
Hence the global top-$k$ can be found among the $k^2$ candidates $\mathcal{I}_1 \times \mathcal{I}_2$.
The same holds for shortlists of any length $k' \ge k$.
\end{proposition}
\begin{proof}
Let $(i,j)$ be among the $k$ largest scores and suppose $i \notin \mathcal{I}_1$. By definition of
$\mathcal{I}_1$ there are $k$ indices $i' \in \mathcal{I}_1$ with $a_{i'} > a_i$ (distinct scores
imply distinct $a$ values within a column). For each of them,
$s_{i'j} = a_{i'} + b_j > a_i + b_j = s_{ij}$. Thus at least $k$ pairs score strictly higher than
$(i,j)$, so $(i,j)$ is not among the $k$ largest --- a contradiction. The argument for $j$ is
symmetric, and a longer shortlist only contains more indices.
\end{proof}

The proposition fails when the shortlist is shorter than $k$: a pair $(i,j)$ with a mediocre $a_i$
and an excellent $b_j$ can be in the global top-$k$ even though $i$ is not in a top-$k'$ list with
$k' < k$. Our numerical check below confirms both directions.

\paragraph{Read and cost.} The selected scores are normalized, $w = \softmax$ over the $k$
winners, and the output is $\vr = \sum_{(i,j)} w_{ij} \vv_{ij}$ with values stored in a large
table $\mV \in \R^{N \times d_v}$. The cost per query is
\begin{equation}
  \underbrace{2\sqrt{N}\cdot \tfrac{d_k}{2}}_{\text{sub-key scores}}
  + \underbrace{k^2}_{\text{candidate sums}}
  + \underbrace{k\, d_v}_{\text{value read}}
  = \bigO\bigl(\sqrt{N}\, d_k + k^2 + k\, d_v\bigr),
\end{equation}
plus two partial sorts of length $\sqrt{N}$ and one of length $k^2$, against $\bigO(N d_k)$ for a
flat search. For $N = 2^{20}$ slots this is $2^{11} = 2048$ half-length dot products instead of a
million full-length ones. Only $k$ rows of $\mV$ receive gradient for a given query, so the
parameter count $N d_v$ can be very large while the update per token stays small.
\citet{c10-lample2019pkm} also use several independent query/sub-key sets (``heads'') and sum
their outputs. Product keys reappear in later sparse-memory and expert-retrieval designs such as
PEER \citep{c10-he2024peer}, memory layers at scale \citep{c10-berges2025memorylayers} and
UltraMem \citep{c10-huang2025ultramem}, in the product-key memories of Hydra
\citep{chaudhary2025hydra} (\cref{ch:hybrid}), and as the first retrieval stage of ExplicitLM
(\cref{sec:c10-explicitlm}).

\begin{example}[$N = 9$, $k = 2$]\label{ex:c10-pkm}
Let $\va = (0.9, 0.1, 0.5)$ and $\vb = (0.2, 0.8, 0.3)$. Then $\mathcal{I}_1 = \{1, 3\}$ and
$\mathcal{I}_2 = \{2, 3\}$, and the four candidates score $s_{12} = 1.7$, $s_{13} = 1.2$,
$s_{32} = 1.3$, $s_{33} = 0.8$. The top two are $(1,2)$ and $(3,2)$. Brute force over all nine
sums ($1.1, 1.7, 1.2, 0.3, 0.9, 0.4, 0.7, 1.3, 0.8$) gives the same two pairs.
\end{example}

\begin{figure}[t]
\centering
\begin{tikzpicture}[x=0.62cm,y=0.62cm,font=\small,
  cell/.style={draw,gray!60,minimum size=0.62cm,inner sep=0pt},
  box/.style={draw,rounded corners=2pt,inner sep=3pt,align=center}]
  \foreach \i in {0,...,5} \foreach \j in {0,...,5} {\node[cell] at (\j,\i) {};}
  \foreach \i/\j in {5/1,5/4,2/1,2/4} {\node[cell,fill=orange!25] at (\j,\i) {};}
  \node[cell,fill=orange!70] at (1,5) {}; \node[cell,fill=orange!70] at (1,2) {};
  \node[box] (q1) at (-5.2,4.0) {$\vq^{(1)}$ vs.\ sub-keys $\bm{c}^{(1)}_i$\\(rows, scores $a_i$)};
  \node[box] (q2) at (-5.2,0.6) {$\vq^{(2)}$ vs.\ sub-keys $\bm{c}^{(2)}_j$\\(columns, scores $b_j$)};
  \draw[-{Stealth}] (q1.east) -- (-0.55,4.0);
  \draw[-{Stealth}] (q2.east) to[out=0,in=-90] (2.5,-0.55);
  \node[box,anchor=west] at (6.2,3.6) {$n \times n = N$ slots\\ $s_{ij} = a_i + b_j$};
  \node[box,anchor=west,fill=orange!10] at (6.2,1.2) {light: $k^2$ candidates\\ dark: final top-$k$};
\end{tikzpicture}
\caption{Product-key lookup. Each half of the query is scored against $\sqrt{N}$ sub-keys; only the
rows and columns of the $k$ best sub-keys (here $k=2$) are combined, and the global top-$k$ is
guaranteed to lie in that $k \times k$ sub-grid.}
\label{fig:c10-pkm}
\end{figure}

\begin{lstlisting}[caption={Reference implementation (ours): product-key top-$k$ retrieval and read.},label={lst:c10-pkm}]
import torch

def product_key_read(q, C1, C2, V, k, k_short=None):
    """q: (B, dk) queries; C1, C2: (n, dk/2) sub-keys; V: (n*n, dv) values.
    Returns the weighted read (B, dv) and the selected slot ids (B, k)."""
    n, h = C1.shape[0], q.shape[-1] // 2
    k_short = k if k_short is None else k_short      # exact only if k_short >= k
    a = q[:, :h] @ C1.T                              # (B, n) first-half scores
    b = q[:, h:] @ C2.T                              # (B, n) second-half scores
    a_val, a_idx = a.topk(k_short, dim=-1)           # (B, k') short list 1
    b_val, b_idx = b.topk(k_short, dim=-1)           # (B, k') short list 2
    cand = a_val[:, :, None] + b_val[:, None, :]     # (B, k', k') outer sum
    cand_id = a_idx[:, :, None] * n + b_idx[:, None, :]
    top_val, top_pos = cand.flatten(1).topk(k, dim=-1)       # (B, k)
    slot = cand_id.flatten(1).gather(1, top_pos)             # (B, k) ids in [0, n*n)
    w = torch.softmax(top_val, dim=-1)                       # (B, k)
    return torch.einsum("bk,bkd->bd", w, V[slot]), slot
\end{lstlisting}

\paragraph{Numerical check.} In our scratch space we compared this procedure (in NumPy) with brute
force over the full grid in 2000 random problems with $n$ between 2 and 39, query widths up to 16
and $k$ up to 16: the selected index sets and scores agreed every time, also with longer
shortlists $k' \ge k$ as ExplicitLM uses. With $n = 30$, $k = 8$ and a shortlist of $k' = 3$
(nine candidates), the result differed from brute force in 1929 of 2000 problems, which
illustrates why $k' \ge k$ is required.

\subsection{Hashing and \texorpdfstring{$N$}{N}-gram embeddings}\label{sec:c10-hashing}

The input embedding table $\bm{E} \in \R^{\abs{\mathcal{V}} \times d}$ is the simplest lookup memory:
the address is the token id $x_t$, the read is the row $\bm{E}[x_t]$, the cost is $\bigO(d)$ and does
not depend on $\abs{\mathcal{V}}$. An \emph{$n$-gram embedding} extends the key to the suffix
$g_{t,n} = (x_{t-n+1},\dots,x_t)$. The key space now has $\abs{\mathcal{V}}^n$ elements --- about
$1.6 \times 10^{10}$ bigrams for a 128k vocabulary --- so a table with one row per possible key is
impossible. \emph{Hashing} solves this: a function $\varphi : \mathcal{V}^n \to \{0,\dots,P-1\}$
maps every $n$-gram to one of $P$ rows of a table $\bm{E}_n \in \R^{P \times d}$. Feature hashing
\citep{c10-weinberger2009hashing}, hash embeddings \citep{c10-svenstrup2017hash} and the hashed
character $n$-grams of fastText \citep{c10-bojanowski2017fasttext} all use this idea, and recent
LLM work hashes token $n$-grams to scale input embeddings
\citep{c10-huang2025overtokenized,c10-yu2025scone,c10-pagnoni2025blt}.

\paragraph{Collisions.} Distinct $n$-grams that hash to the same row share their embedding.
Model $\varphi$ as a uniformly random function (the guarantee that universal hash families give
for pairs; \citealp{c10-carter1979universal}). Let $U$ distinct $n$-grams occur in the data. For a
fixed $n$-gram, each of the other $U-1$ lands in its row independently with probability $1/P$,
so
\begin{equation}
  \Pr[\text{the $n$-gram shares its row}] = 1 - \Bigl(1 - \frac{1}{P}\Bigr)^{U-1}
  \approx 1 - e^{-U/P},
  \qquad
  \E[\#\text{colliding pairs}] = \binom{U}{2}\frac{1}{P} .
  \label{eq:c10-collide}
\end{equation}
With as many rows as keys ($U = P$), about $1 - e^{-1} \approx 63\%$ of the $n$-grams share a row
with some other $n$-gram. Collisions are therefore not rare events to be ignored; they are the
normal state of a hashed table, and the model must be able to cope with them.

\paragraph{Multi-head hashing.} Use $K$ independent hash functions $\varphi_1,\dots,\varphi_K$
into $K$ tables of sizes $P_1,\dots,P_K$ and concatenate (or sum) the $K$ rows. Two $n$-grams are
now indistinguishable only if they collide in \emph{every} head, which under the random-function
model has probability $\prod_j 1/P_j$, so that even with $U \approx P_j$ almost every $n$-gram
gets a unique \emph{combination} of rows although each individual row is shared.

\paragraph{A multiplicative--XOR hash.} Engram (\cref{sec:c10-engram}) uses a cheap deterministic
hash. With odd random multipliers $m_0, m_1, \dots$ (one per position in the $n$-gram), it forms
\begin{equation}
  \mathrm{mix}(g_{t,n}) = (x_t m_0) \oplus (x_{t-1} m_1) \oplus \dots \oplus (x_{t-n+1} m_{n-1}),
  \qquad
  z_{t,n,j} = \mathrm{mix}(g_{t,n}) \bmod P_{n,j},
  \label{eq:c10-mixhash}
\end{equation}
where $\oplus$ is bitwise exclusive-or and the table sizes $P_{n,j}$ are distinct primes. Different
multipliers per position make the hash order-sensitive ($(a,b)$ and $(b,a)$ mix differently);
multiplying by a large odd number spreads the bits of a small token id over the whole word; the
prime modulus spreads the result over the table. The same mixed value is reduced modulo $K$
different primes, which has a consequence worth stating precisely.

\begin{proposition}[All-head collisions reduce to collisions of the mix]\label{prop:c10-crt}
Let $P_1,\dots,P_K$ be distinct primes and let $a, b$ be integers with $0 \le a, b < \prod_j P_j$.
Then $a \bmod P_j = b \bmod P_j$ for all $j$ if and only if $a = b$.
\end{proposition}
\begin{proof}
If $a \equiv b \pmod{P_j}$ for every $j$, each $P_j$ divides $a - b$. Distinct primes are pairwise
coprime, so their product divides $a-b$ (this is the uniqueness part of the Chinese remainder
theorem). Since $\abs{a - b} < \prod_j P_j$, the only multiple available is $a - b = 0$. The
converse is trivial.
\end{proof}

In Engram's demo configuration the mixed values are non-negative 64-bit integers below $2^{63}$,
while eight primes near $6.5 \times 10^5$ multiply to about $3 \times 10^{46}$. By
\cref{prop:c10-crt}, two $n$-grams receive identical rows in all eight heads exactly when their
mixed values coincide. The heads therefore do not reduce the rate of \emph{complete} collisions
below that of the 63-bit mix (which is tiny anyway); what they buy is that the \emph{partial}
collisions, where some rows are shared, are spread over different partners in different heads.

\paragraph{Numerical check.} We re-implemented the hash of \cref{eq:c10-mixhash} in NumPy, both
vectorized (shifting the token array, as the Engram demo does) and as a naive loop over positions,
and found identical indices. On 200{,}000 random pairs of 62-bit values, ``same index in all eight
heads'' coincided exactly with ``same value''; the single-head collision rate for a prime of 1009
was $9.95 \times 10^{-4}$ against the predicted $1/1009 = 9.91 \times 10^{-4}$. With the small
primes $3, 5, 7$ (product 105), values that differ by multiples of 105 collided in all heads, as
the proposition predicts when its range condition fails.

\paragraph{Why a lookup is $\bigO(1)$.} Computing $z_{t,n,j}$ costs $\bigO(n)$ integer operations
and reading the row costs $\bigO(d)$ memory traffic. Neither depends on $P$, so the table can grow
until it no longer fits in memory without changing the per-token compute. This is the property
that makes hashed tables attractive as a place to put parameters.

\subsection{Mixture-of-Experts as conditional parameter memory}\label{sec:c10-moe}

The MoE layer (\cref{ch:foundations}) replaces one FFN by $E$ expert FFNs and a router. Given a
hidden state $\vh$, the router computes logits $\vr = \mW_r \vh \in \R^{E}$, selects a small set
$\mathcal{S}(\vh) = \TopK(\vr, k)$, and the layer outputs, in the survey's notation
(\S IV-C, Eq.~(1)),
\begin{equation}
  \vy = \sum_{i \in \mathcal{S}(\vh)} g_i(\vh)\, \mathrm{Expert}_i(\vh),
  \label{eq:c10-moe}
\end{equation}
with gating weights $g_i(\vh)$ such as a softmax over the selected logits
\citep{c01-shazeer2017moe}. Switch Transformer \citep{fedus2022switch} routes each token to a
single expert, GLaM \citep{du2022glam} scales the idea to very large total parameter counts,
Mixtral \citep{jiang2024mixtral} uses top-2 routing over eight experts per layer, and DeepSeekMoE
\citep{dai2024deepseekmoe} splits experts into many fine-grained ones and adds always-active shared
experts.

From the memory point of view, each expert is a persistent block of parameters and the router is
a learned content-based address over $E$ ``slots'', each slot a whole function. Two parameter
counts describe the trade-off. If every expert has $P_e$ parameters, the layer stores
$P_{\mathrm{tot}} = E P_e$ parameters but activates only $P_{\mathrm{act}} = k P_e$ per token, so
the per-token compute is that of a $k P_e$-parameter layer while the capacity is that of an
$E P_e$-parameter layer. The survey therefore places MoE between parameterized and lookup memory:
it ``transforms parameter space into an addressable memory system governed by learned routing''
(\S IV-C). It is offline and long-term in the taxonomy.

The contrast with a hashed table is the starting point of Engram. An MoE address depends on the
hidden state $\vh$, so it is known only when the layer runs, and a selected slot is \emph{computed}
(an MLP evaluation costing about $2 P_e$ floating-point operations). A hashed address depends only
on token ids, so it is known before the forward pass, and a selected slot is \emph{copied} (a row
read). Engram calls the first conditional computation and the second conditional memory.

\subsection{Precursors: PlugLM and MEMORYLLM}\label{sec:c10-precursors}

Two pre-2025 systems anticipate the designs of this chapter.

\textbf{PlugLM} \citep{cheng2023pluglm} starts from \cref{eq:c10-ffn-kv}: if FFNs already behave
like key--value stores, replace selected FFN layers by a differentiable plug-in key--value memory
whose entries are retrieved by a ``knowledge attention'' step and can be updated without retraining
the model (survey \S IV-B(d)). It sits between parameter memory and external lookup and is the
direct ancestor of ExplicitLM's goal of editable knowledge.

\textbf{MEMORYLLM} \citep{wang2024memoryllm} (2024; ``Towards Self-Updatable Large Language
Models'') embeds a fixed-size pool of latent memory tokens in every Transformer layer and refreshes
it with forward-only self-updates when new knowledge is read, rewriting dedicated memory tokens
while the backbone stays unchanged (survey \S IV-A(b), (e)). The survey lists it as online and
long-term. It must not be confused with the 2026 \emph{MemoryLLM} of \cref{sec:c10-memoryllm},
which is an offline design that turns FFNs into token-indexed lookups.

With these tools in hand, the four 2025--2026 systems can be read as different answers to the
same questions: what is stored (slot vectors, $N$-gram embeddings, token sequences, FFN outputs),
how it is addressed (attention, hashing, product keys, token id), and when it is written.

% =====================================================================================
\section{LM2: a gated memory bank in every decoder layer}\label{sec:c10-lm2}

\subsection{Motivation}

Large Memory Models (LM2) \citep{kang2025lm2} target a weakness that long-context benchmarks
expose: standard Transformers struggle with multi-step reasoning, relational argumentation and the
synthesis of facts that are spread across a long context. Two earlier families of memory-augmented
models address this with a separate set of memory vectors. The Recurrent Memory Transformer
\citep{c10-bulatov2022rmt} passes memory tokens from one segment to the next through the ordinary
attention layers, and relational memory cores \citep{c10-santoro2018relational} update a matrix
of memory slots with attention and LSTM-style gates \citep{hochreiter1997lstm}. LM2's own
\code{MemoryModule} describes itself as a ``relational memory core with multi-head attention and
gating''.

LM2's design choice is to leave the Transformer's normal information flow untouched and add a
\emph{complementary memory pathway} in every decoder layer: a bank of memory slots that the
tokens read through cross-attention, whose contribution to the residual stream is gated, and that
is rewritten through input and forget gates. Unlike retrieval systems the slots are trained as
part of the model; unlike an implicit recurrent state they are designated memory variables with a
separate storage role (survey \S IV-A(d)).

\begin{taxonomy}
\textbf{Representation:} explicit --- dedicated memory slots with their own read (cross-attention)
and write (gates) interface \citep[\S IV-A(d)]{zhoubian2026memory}.
\textbf{Update dynamics:} offline according to the survey's Table~I (the slot parameters and gates
are learned in training). In the official code the bank is also rewritten by the gates in every
forward pass and carried from one batch to the next, so it behaves like an online state during
training.
\textbf{Persistence:} long-term (Table~I).
\textbf{Update rule:} gated slot writes; the survey groups LM2 with MEMORYLLM as ``slot-oriented
or forward-update variants'' whose memory is structurally separated from the backbone
(\S IV-A(e)). Table~I's paradigm label is ``Dedicated Memory Slots''.
\end{taxonomy}

\subsection{Formulation}

\paragraph{Notation.} To avoid a clash with the causal mask $\mM$, we call the memory bank $\mB$.
It has $N$ slots of width $d$, $\mB \in \R^{N \times d}$, with row $\vb_r$ for slot $r$. The
token representations that enter the memory module of a layer are the rows of
$\bm{E} \in \R^{T \times d}$ (in the code, the output of the layer's self-attention). The memory
module has its own projections $\mU_Q, \mU_K, \mU_V, \mU_{\mathrm{out}}, \mU_{\mathrm{in}},
\mU_{\mathrm{f}} \in \R^{d \times d}$. The equations below are our transcription of the paper's
description, with the shapes made explicit; consult the paper for its exact parameterization, and
see \cref{sec:c10-lm2-code} for what the code does.

\paragraph{Initialization.} Every slot starts as a basis vector: $\mB^{(0)} = \mI$, which needs
$N = d$. The slots are then orthonormal and equally distant from one another, so the first reads
are not degenerate (with all-zero or identical slots every read would return the same vector).

\paragraph{Read by cross-attention.} The tokens are the queries and the slots provide keys and
values:
\begin{equation}
  \mQ = \bm{E}\mU_Q,\quad \mK = \mB\mU_K,\quad \mV = \mB\mU_V,\qquad
  \bm{E}_{\mathrm{mem}} = \softmax\Bigl(\frac{\mQ\mK\T}{\sqrt{d}} + \mM_{\mathrm{mem}}\Bigr)\mV
  \in \R^{T\times d},
  \label{eq:c10-lm2-read}
\end{equation}
where $\mM_{\mathrm{mem}} \in \R^{T \times N}$ is a mask. The paper applies causal masking; the code
uses $(\mM_{\mathrm{mem}})_{tr} = -\infty$ for $r > t$, so token $t$ may read slots $1,\dots,t$.
Row $t$ of $\bm{E}_{\mathrm{mem}}$ is the memory read-out for token $t$.

\paragraph{Gated output.} An output gate decides how much memory flows into the residual stream:
\begin{equation}
  \mG_{\mathrm{out}} = \sigmoid\bigl(\bm{E}_{\mathrm{mem}}\mU_{\mathrm{out}}\bigr),\qquad
  \bm{E}_{\mathrm{gated}} = \mG_{\mathrm{out}} \had \mB ,
  \label{eq:c10-lm2-out}
\end{equation}
and the decoder layer adds it next to the self-attention output, before the FFN:
\begin{equation}
  \mX' = \mX + \mathrm{Attn}\bigl(\LayerNorm(\mX)\bigr) + \bm{E}_{\mathrm{gated}}, \qquad
  \mX_{\mathrm{next}} = \mX' + \mathrm{FFN}\bigl(\LayerNorm(\mX')\bigr).
  \label{eq:c10-lm2-block}
\end{equation}
(The Llama backbone of the code uses RMSNorm; we write $\LayerNorm$ generically.) Note the shapes:
$\mG_{\mathrm{out}}$ is $T \times d$ and $\mB$ is $N \times d$, so the elementwise product aligns
token $t$ with slot $t$ and requires $N = T$. The official code enforces exactly this.

\paragraph{Gated write.} Input and forget gates produce the next bank, in the form of an LSTM cell
update with the read-out as the candidate:
\begin{equation}
  \mG_{\mathrm{in}} = \sigmoid\bigl(\bm{E}\mU_{\mathrm{in}}\bigr),\quad
  \mG_{\mathrm{f}} = \sigmoid\bigl(\bm{E}_{\mathrm{mem}}\mU_{\mathrm{f}}\bigr),\qquad
  \mB_{\mathrm{next}} = \mG_{\mathrm{in}} \had \tanh\bigl(\bm{E}_{\mathrm{mem}}\bigr)
  + \mG_{\mathrm{f}} \had \mB .
  \label{eq:c10-lm2-write}
\end{equation}
The comparison with an LSTM cell $\vs_t = \vg^{f}_t \had \vs_{t-1} + \vg^{i}_t \had \tanh(\cdot)$
explains the roles: the forget gate keeps old slot content, the input gate admits new content, and
$\tanh$ bounds each new contribution to $(-1,1)$.

\begin{example}[One slot coordinate]
Take $b = 1$ (an identity-initialized slot on its own axis), a read-out coordinate $0.4$, an input
gate $0.5$ and a forget gate $\sigmoid(1) \approx 0.731$ (the code's forget-gate bias is
initialized to $1$). Then $b_{\mathrm{next}} = 0.5\tanh(0.4) + 0.731 \cdot 1 = 0.5 \cdot 0.380 +
0.731 \approx 0.921$: the slot decays slightly towards the new content.
\end{example}

\paragraph{How large can the bank grow?} Each entry is updated by
$b' = g_{\mathrm{in}}\tanh(e) + g_{\mathrm{f}}\, b$ with $g_{\mathrm{in}}, g_{\mathrm{f}} \in (0,1)$.
\begin{proposition}[Growth of the bank]\label{prop:c10-lm2-bound}
Let $\beta_s = \max_{r,c} \abs{(\mB^{(s)})_{rc}}$ after $s$ writes. Then $\beta_{s+1} \le 1 +
\beta_s$, so $\beta_s \le \beta_0 + s$. If in addition every forget gate satisfies
$g_{\mathrm{f}} \le \gamma < 1$, then $\beta_s \le \max\bigl(\beta_0, 1/(1-\gamma)\bigr)$ for all $s$.
\end{proposition}
\begin{proof}
$\abs{b'} \le g_{\mathrm{in}}\abs{\tanh e} + g_{\mathrm{f}}\abs{b} \le 1 + g_{\mathrm{f}}\,\beta_s$.
With $g_{\mathrm{f}} < 1$ this gives the first bound. With $g_{\mathrm{f}} \le \gamma$, the map
$\beta \mapsto 1 + \gamma\beta$ has fixed point $\beta^\star = 1/(1-\gamma)$; if $\beta_s \le
\max(\beta_0,\beta^\star)$ then $\beta_{s+1} \le 1 + \gamma\max(\beta_0,\beta^\star) \le
\max(\beta_0,\beta^\star)$, because $1 + \gamma\beta \le \beta$ whenever $\beta \ge \beta^\star$.
Induction completes the proof.
\end{proof}
With the initial forget bias, $\gamma \approx 0.73$ gives $\beta^\star \approx 3.7$. Nothing in the
architecture keeps the forget gate away from $1$, however, so over many writes (the bank is written
$L$ times per forward pass and carried across batches) the only guaranteed bound is linear.

\paragraph{Two clocks.} The bank is a recurrent state, but ``time'' has two meanings. Within one
forward pass the bank is threaded through the depth of the network:
$\mB^{(\ell)} = \mathrm{Write}_\ell\bigl(\mB^{(\ell-1)}, \bm{E}^{(\ell)}\bigr)$ for
$\ell = 1,\dots,L$, so each layer reads what earlier layers wrote. Across forward passes, the code
keeps the final bank $\mB^{(L)}$ of one batch as the initial bank of the next, with the gradient
cut (\code{memory = self.memory.detach()}). This is the recurrent-memory-across-segments regime of
RMT, trained with truncation after one segment.

\paragraph{Parallel versus per-token view.} \Cref{eq:c10-lm2-read,eq:c10-lm2-out,eq:c10-lm2-write}
process all $T$ tokens at once with a mask. Row $t$ of every output only involves row $t$ of
$\bm{E}$, rows $1,\dots,t$ of $\mB$ (through the masked softmax) and row $t$ of $\mB$ (through the
gates), so the same results come out of a loop over $t$ that reads only those rows. Our NumPy check
below confirms that the masked parallel form and the loop agree to machine precision. The
per-token view also gives the causality argument.

\begin{proposition}[Causality of the per-token-gated layer]\label{prop:c10-lm2-causal}
Assume the bank entering the first layer does not depend on the current segment. With the gates of
\cref{eq:c10-lm2-write} computed row by row and the slot mask $(\mM_{\mathrm{mem}})_{tr} = -\infty$
for $r > t$, row $t$ of every layer's output and rows $1,\dots,t$ of every layer's bank depend only
on tokens $1,\dots,t$.
\end{proposition}
\begin{proof}
By induction over layers. The claim holds for the input bank. Suppose rows $1,\dots,t$ of
$\mB^{(\ell-1)}$ and of the layer input depend only on tokens $\le t$. Self-attention is causal.
Row $t$ of $\bm{E}_{\mathrm{mem}}$ is computed from row $t$ of $\bm{E}$ and rows $\le t$ of
$\mB^{(\ell-1)}$; rows $t$ of $\mG_{\mathrm{out}}, \mG_{\mathrm{in}}, \mG_{\mathrm{f}}$ are functions
of rows $t$ of $\bm{E}$ and $\bm{E}_{\mathrm{mem}}$; row $t$ of the output and of
$\mB^{(\ell)}$ combine these with row $t$ of $\mB^{(\ell-1)}$. All depend on tokens $\le t$ only.
\end{proof}

\begin{pitfall}[title={Pitfall: pooling over the sequence breaks causality}]
A gate computed from a statistic of the whole sequence --- for example the mean of the layer input
over all $T$ positions --- makes row $t$ of the bank, and therefore the output at position $t$,
depend on tokens after $t$. In next-token training with teacher forcing this leaks the targets
into the inputs. The official LM2 code computes its input-gate contribution from such a mean
(\cref{sec:c10-lm2-code}). In a two-layer NumPy transcription with $T = N = d = 6$, perturbing only
the last token changed the outputs at earlier positions by up to $0.18$; with per-token gates the
change was exactly zero. The same issue appears in the query of the ExplicitLM code
(\cref{sec:c10-explicitlm}).
\end{pitfall}

\subsection{Algorithms and complexity}

\begin{algorithm}[t]
\caption{One LM2 decoder layer (parallel form, per-token gates)}\label{alg:c10-lm2-layer}
\begin{algorithmic}[1]
\Require layer input $\mX \in \R^{T\times d}$, bank $\mB \in \R^{N\times d}$ with $N = T$
\Ensure layer output $\mX_{\mathrm{next}}$, next bank $\mB_{\mathrm{next}}$
\State $\mA \gets \mathrm{Attn}(\LayerNorm(\mX))$ \Comment{causal self-attention, unchanged}
\State $\bm{E} \gets \mA$ \Comment{tokens that query the memory (as in the code)}
\State $\mQ \gets \bm{E}\mU_Q$;\; $\mK \gets \mB\mU_K$;\; $\mV \gets \mB\mU_V$
\State $\bm{E}_{\mathrm{mem}} \gets \softmax(\mQ\mK\T/\sqrt{d} + \mM_{\mathrm{mem}})\mV$
\State $\mX' \gets \mX + \mA + \sigmoid(\bm{E}_{\mathrm{mem}}\mU_{\mathrm{out}}) \had \mB$
\State $\mX_{\mathrm{next}} \gets \mX' + \mathrm{FFN}(\LayerNorm(\mX'))$
\State $\mB_{\mathrm{next}} \gets \sigmoid(\bm{E}\mU_{\mathrm{in}}) \had \tanh(\bm{E}_{\mathrm{mem}})
       + \sigmoid(\bm{E}_{\mathrm{mem}}\mU_{\mathrm{f}}) \had \mB$
\State \Return $\mX_{\mathrm{next}}, \mB_{\mathrm{next}}$
\end{algorithmic}
\end{algorithm}

\begin{algorithm}[t]
\caption{LM2 training on a stream of segments (as organized in the official code)}\label{alg:c10-lm2-train}
\begin{algorithmic}[1]
\Require segments $\mX_1, \mX_2, \dots$ of length $T$, number of layers $L$
\State $\mB \gets \mI_N$ \Comment{one identity bank per sequence in the batch}
\For{$s = 1, 2, \dots$}
  \State $\mB^{(0)} \gets \mathrm{stopgrad}(\mB)$;\; $\mH^{(0)} \gets \mathrm{Embed}(\mX_s)$
  \For{$\ell = 1,\dots,L$}
    \State $(\mH^{(\ell)}, \mB^{(\ell)}) \gets \mathrm{Layer}_\ell(\mH^{(\ell-1)}, \mB^{(\ell-1)})$
    \Comment{\cref{alg:c10-lm2-layer}}
  \EndFor
  \State $\loss \gets$ next-token cross-entropy of $\mH^{(L)}$;\; update all weights by AdamW
  \State $\mB \gets \mB^{(L)}$ \Comment{carried to the next segment, gradient cut}
\EndFor
\end{algorithmic}
\end{algorithm}

\paragraph{Cost.} Per layer, the projections cost $\bigO(T d^2 + N d^2)$ and the read
\eqref{eq:c10-lm2-read} costs $\bigO(T N d)$ time and $\bigO(T N)$ memory for the attention weights.
With $N = T$ the read is as expensive as self-attention itself, so LM2 roughly doubles the
attention-like work of every layer. The state passed between layers and segments is the bank,
$N d$ numbers per sequence. For autoregressive decoding the mirrored code has no incremental path
(the sequence length must equal $N$). One natural option, which is our suggestion and not part of
the paper or code, is to read a fixed bank during a segment at $\bigO(N d)$ per token and commit the
gated write at the segment boundary.

\subsection{Implementation walkthrough}\label{sec:c10-lm2-code}

The official repository \repo{convergence-ai__lm2} (license CC BY-NC 4.0) is a compact PyTorch
project. The memory module is in \repofile{repos/convergence-ai__lm2/src/memory.py};
the file \repofile{repos/convergence-ai__lm2/src/model_memory_llama.py} wires it into the Hugging
Face Llama model;
\repofile{repos/convergence-ai__lm2/train.py} and
\repofile{repos/convergence-ai__lm2/src/trainer.py} run distributed training, and
\repofile{repos/convergence-ai__lm2/data_proc/smollm.py} prepares the data.

\begin{implnote}[title={In the code: wiring into each decoder layer}]
\code{LlamaMem} subclasses \code{LlamaForCausalLM} and replaces every decoder layer by a
\code{CustomLlamaDecoderLayer}, so each layer has its own \code{MemoryAttention} and its own
\code{MemoryModule} parameters. \code{MemoryAttention.forward} first runs the layer's ordinary
self-attention, then passes the attention \emph{output} (not the layer input) to the memory module
as $\bm{E}$. The module returns two tensors: the attention output plus the new bank, which goes to
the residual stream, and the new bank, which is passed to the next layer
(\cref{lst:c10-lm2-layer}). \code{LlamaMem.forward} threads the bank through the layer loop and
stores the final one in \code{self.memory} (and in the buffer \code{memory_bank}, saved with
checkpoints), so the next call starts from it.
\end{implnote}

\begin{lstlisting}[caption={Excerpt from \texttt{repos/convergence-ai\_\_lm2/src/model\_memory\_llama.py} (\texttt{CustomLlamaDecoderLayer.forward}).},label={lst:c10-lm2-layer}]
        residual = hidden_states
        hidden_states = self.input_layernorm(hidden_states)

        if self.use_memory and memory is not None:
            attn_output, memory = self.mem_attn(
                hidden_states=hidden_states,
                memory=memory,
                attention_mask=attention_mask,
                position_ids=position_ids,
                position_embeddings=position_embeddings,
            )
        else:
            attn_output, _, _ = self.self_attn(
                hidden_states,
                attention_mask,
                position_ids=position_ids,
                position_embeddings=position_embeddings,
            )

        hidden_states = residual + attn_output
        residual = hidden_states
        hidden_states = self.post_attention_layernorm(hidden_states)
        hidden_states = self.mlp(hidden_states)
        hidden_states = residual + hidden_states

        return hidden_states, memory
\end{lstlisting}

\begin{lstlisting}[caption={Excerpt from \texttt{repos/convergence-ai\_\_lm2/src/memory.py} (\texttt{MemoryModule.forward}).},label={lst:c10-lm2-mem}]
        inputs = inputs.view(inputs.shape[0], inputs.shape[1], -1)

        # Attend over memory to obtain the next memory state
        next_memory = self.attend_over_memory(inputs, memory, attention_mask)

        # Create input and forget gates based on current inputs and memory
        input_gate, forget_gate = self.create_gates(inputs, memory)
        # Apply input gate with tanh activation
        next_memory = input_gate * torch.tanh(next_memory)
        # Apply forget gate to retain part of the previous memory
        next_memory += forget_gate * memory

        # Return the updated inputs and the next memory state
        return inputs + next_memory, next_memory
\end{lstlisting}

\begin{implnote}[title={In the code: shapes and the $N = T = d$ constraint}]
The bank is created as \code{torch.stack([torch.eye(memory_slots) ...])}, one $M \times M$ identity
per sequence in the batch, where $M$ is \code{memory_slots}. Inside \code{multi_head_attention},
\code{query_proj} maps the width-$d$ tokens to width $M$, \code{key_proj} and \code{value_proj}
are $M \times M$, the scores are scaled by $1/\sqrt{M/H_{\mathrm{mem}}}$ for
$H_{\mathrm{mem}}$ = \code{num_mem_heads} heads, and a boolean mask \code{~torch.tril(...)} blocks
slot $r > t$ for token $t$. The code raises an error unless the sequence length equals $M$, and
\code{inputs + next_memory} adds a $(B, M, M)$ bank to $(B, T, d)$ token states, so in practice
$M = T = d$. The configuration in \repofile{repos/convergence-ai__lm2/configs/model/llama3.2-1b.yaml}
sets \code{memory_slots: 2048} and \code{num_mem_heads: 8}, and
\repofile{repos/convergence-ai__lm2/scripts/train.sh} sets \code{sequence_length=2048}; the model
width comes from the Llama-3.2-1B configuration, which must therefore also be 2048.
\end{implnote}

\begin{implnote}[title={In the code: the gates differ from the equations}]
\code{create_gates(inputs, memory)} computes both gates from one sum. The input part is
\code{RepeatLinear}: a learnable vector $\vw$ scales the tokens, $\ReLU(\vw \had \ve_t)$ is
\emph{averaged over all $T$ positions}, and a linear map produces $2M$ gate logits shared by all
slots. The memory part is a linear map of $\tanh(\mB)$ (\code{GroupLinearLayer}, uniform
initialization), giving $2M$ logits per slot. Their sum is split into input and forget logits, and
biases are added before the sigmoid (\code{input_bias} initialized to 0, \code{forget_bias} to 1).
The candidate is not the raw read-out: \code{attend_over_memory} forms
$\LayerNorm(\mB + \bm{E}_{\mathrm{mem}})$, applies a two-layer ReLU MLP with a residual connection
and a second LayerNorm, and only then the gates are applied
(\cref{lst:c10-lm2-mem}). Compared with \cref{eq:c10-lm2-out,eq:c10-lm2-write}: (i) there is no
separate output gate --- the \emph{gated new bank} is added to the residual stream; (ii) the input
gate depends on a sequence-wide mean, which is not causal (see the pitfall above); (iii) the MLP is
built as \code{nn.ModuleList([nn.Linear(M, M)] * 2)}, which repeats the \emph{same} module twice,
so its two layers share weights; (iv) the attention dropout defaults to \code{attn_drop=0.9} and
\code{MemoryAttention} does not override it, so 90\% of the memory attention weights are dropped
in training mode. Whether (ii)--(iv) are intended we cannot tell from the repository.
\end{implnote}

\begin{implnote}[title={In the code: training setup}]
\code{train.py} reads Hydra configs (\repofile{repos/convergence-ai__lm2/configs/train.yaml} and
its \code{model}, \code{train}, \code{pretrain} groups). The architecture is that of
\code{AutoConfig.from_pretrained("meta-llama/Llama-3.2-1B")} extended by \code{CustomLlamaConfig}
with \code{use_memory}, \code{memory_slots} and \code{num_mem_heads}; the model is built with
\code{LlamaMem.from_config}, i.e.\ \emph{trained from scratch}, not fine-tuned from Llama weights
(the YAML fields \code{n_layer}, \code{n_embd} are copied into an unused config object).
\code{scripts/train.sh} launches \code{torchrun} with batch size 1 per GPU, sequence length 2048,
bfloat16 autocast, learning rate $2\times10^{-4}$, 700 warm-up steps, \code{lr_decay_frac=0.0},
600{,}000 iterations and \code{model.use_memory=true}. The trainer uses AdamW with betas
$(0.9, 0.95)$, weight decay $0.1$ on linear weights only, and gradient clipping at $1.0$. The
schedule is linear warm-up followed by cosine decay to
\code{learning_rate * lr_decay_frac}; with \code{lr_decay_frac=0.0} it therefore decays to zero,
although the README describes $0.0$ as ``no decay''. Data are the SmolLM corpus subsets
(\code{cosmopedia-v2} by default) tokenized with the Llama-3 tokenizer into binary shards with a
256-integer header (magic number 20240801), read by \code{DistributedDataLoader}. A validation
pass on 100 batches runs every \code{log_freq} steps; note that it also overwrites the carried
bank.
\end{implnote}

\paragraph{Parameter count.} From the module definitions, one \code{MemoryModule} with $M = d$
has $3(M^2+M)$ parameters in the query/key/value projections, $M^2 + M$ in the (shared) MLP,
$2(2M^2 + 2M)$ in the two gate projectors, $M$ in $\vw$, $4M$ in the two LayerNorms and two
biases: $33{,}581{,}058$ parameters for $M = 2048$ (about 33.6M per decoder layer, our count), to
be multiplied by the number of decoder layers of the base configuration.

\begin{lstlisting}[caption={Reference implementation (ours): LM2 memory read, gated output and gated write with per-token gates.},label={lst:c10-lm2-ref}]
import torch, torch.nn as nn

class LM2Memory(nn.Module):
    """Memory pathway of one decoder layer. E: (B, T, d) token states,
    bank: (B, N, d) slots with N == T (slot t is aligned with token t)."""
    def __init__(self, d):
        super().__init__()
        self.q, self.k, self.v = (nn.Linear(d, d, bias=False) for _ in range(3))
        self.w_out = nn.Linear(d, d)
        self.w_in = nn.Linear(d, d)
        self.w_f = nn.Linear(d, d)
        nn.init.constant_(self.w_f.bias, 1.0)          # start by remembering

    def forward(self, E, bank):
        B, T, d = E.shape
        N = bank.shape[1]
        assert N == T, "elementwise gates align slot t with token t"
        q, k, v = self.q(E), self.k(bank), self.v(bank)
        s = q @ k.transpose(1, 2) / d ** 0.5           # (B, T, N)
        future = torch.ones(T, N, dtype=torch.bool, device=E.device).triu(1)
        e_mem = s.masked_fill(future, float("-inf")).softmax(-1) @ v   # (B, T, d)
        g_out = torch.sigmoid(self.w_out(e_mem))       # output gate
        e_gated = g_out * bank                         # added to the residual stream
        g_in = torch.sigmoid(self.w_in(E))             # per-token input gate
        g_f = torch.sigmoid(self.w_f(e_mem))           # per-token forget gate
        new_bank = g_in * torch.tanh(e_mem) + g_f * bank
        return e_gated, new_bank

def init_bank(batch, d):
    return torch.eye(d).expand(batch, d, d).clone()   # N = d identity slots
\end{lstlisting}

\paragraph{Numerical checks.} In NumPy (scratch space) we verified that (a) the masked parallel
form of \cref{lst:c10-lm2-ref} equals a per-token loop that reads only slots $\le t$ (outputs and
new bank agree to machine precision); (b) with per-token gates, perturbing the last token leaves
all earlier outputs of a two-layer stack unchanged, while a transcription of the code's
sequence-pooled gate changes them (by up to $0.18$ in our toy setting); and (c) the parameter
count above.

\begin{results}
\textbf{Source files:} \repofile{repos/convergence-ai__lm2/README.md}; survey PDF.
The repository README reports no numbers, and the survey discusses LM2 only qualitatively. The
paper's abstract reports that LM2 outperforms a memory-augmented baseline (RMT) and a Llama-3.2
baseline on the BABILong long-context reasoning benchmark, with strong multi-hop inference,
numerical reasoning and long-context question answering, and that it improves over a pre-trained
vanilla model on MMLU, i.e.\ the memory pathway does not hurt general ability. See the paper for
the numbers.
\end{results}

\subsection{Discussion}

LM2 shows that explicit memory can be wired in at the level of every layer as a persistent
representational workspace rather than bolted on as an external datastore (survey \S IV-A(d)).
Its strengths are the untouched main pathway (memory is additive and gated, so the model can learn
to ignore it) and a write rule with well-understood dynamics: \cref{eq:c10-lm2-write} is the gated
state transition of an LSTM, and in the language of \cref{ch:linear} it is a recurrence with
data-dependent decay $\mG_{\mathrm{f}}$ and write strength $\mG_{\mathrm{in}}$, applied to slots
instead of a single matrix state.

Its limitations are equally clear. The bank is small: $N d$ numbers per sequence, whereas the
lookup memories below hold billions of parameters; LM2's memory is closer to working memory than
to a knowledge store. Aligning slot $t$ with token $t$ ties $N$ to the sequence length, and the read
is as expensive as attention. In the official code, carrying the bank from one training batch to
the next means that unrelated documents write into the same state, and the gradient never crosses
a batch boundary, so the model cannot learn \emph{what} to store for later segments beyond what
helps the current one. Compared with test-time training (\cref{ch:ttt}), where Titans and TTT
write memory by gradient steps on a self-supervised loss, LM2 writes by learned gates in a single
forward pass; compared with MEMORYLLM (\cref{sec:c10-precursors}), it rewrites all slots softly
instead of replacing a subset.

% =====================================================================================
\section{Engram: conditional memory via scalable lookup}\label{sec:c10-engram}

\subsection{Motivation}

Engram \citep{cheng2026engram} starts from an observation about language. Modeling text involves
two qualitatively different sub-tasks: compositional reasoning, which needs deep and dynamic
computation, and knowledge retrieval, much of which is \emph{local, static and stereotyped} ---
named entities, idioms, formulaic phrases. Classical $N$-gram models capture such local regularities
with cheap lookups, but a Transformer has no lookup primitive and must \emph{simulate} retrieval
with computation. The Engram paper illustrates this with an entity-resolution trace reproduced
from earlier interpretability work: to build the representation of ``Diana, Princess of Wales'' a
model spends about six layers of attention and FFNs composing ``Wales'', ``Princess of Wales'' and
finally the person. That is depth spent on reconstructing a static table at run time.

MoE (\cref{sec:c10-moe}) already scales capacity by \emph{conditional computation}. Engram adds
\emph{conditional memory}: a module, inserted at a few layers, that maps the last few tokens to
rows of very large hashed $N$-gram embedding tables in $\bigO(1)$ time, and lets the current hidden
state decide how much of the retrieved vector to use. The paper formulates how to split a fixed
budget of sparse parameters between MoE experts and Engram tables, and designs the module so that
the tables can live in host memory. The survey highlights it because it separates memory sparsity
from expert sparsity: ``MoE layers sparsely activate parameterized transformations, whereas Engram
sparsely addresses stored memory entries'' (\S IV-B(c)).

\begin{taxonomy}
\textbf{Representation:} explicit --- hashed lookup slots that persist and are individually
addressed \citep[\S IV-B(c)]{zhoubian2026memory}; Table~I's paradigm label is ``Hashed Lookup Slots''.
\textbf{Update dynamics:} offline --- tables are trained with the model and fixed at inference.
\textbf{Persistence:} long-term.
\textbf{Update rule:} objective-induced or structural (Table~II lists Engram there together with MoE
models): memory content is shaped by the training objective and the offline construction of the
tables, not by per-instance online writes.
\end{taxonomy}

\subsection{Formulation}

Let $\vh_t \in \R^{d}$ be the hidden state at position $t$ entering an Engram layer and $x_t$ the
token at position $t$. The module works in two phases, retrieval and fusion
(\cref{fig:c10-engram}).

\paragraph{Tokenizer compression.} Subword tokenizers give different ids to strings that are the
same for the purpose of a lookup (\emph{Apple} and \emph{\textvisiblespace apple}). Engram first
applies a precomputed surjection $\mathcal{P} : \mathcal{V} \to \mathcal{V}'$ that maps each raw id
to a canonical id by normalizing its text (Unicode NFKC, lowercasing, and so on):
$x'_t = \mathcal{P}(x_t)$. The paper reports a 23\% reduction of the effective vocabulary of its
128k tokenizer. Fewer distinct ids means fewer distinct $n$-grams competing for table rows.

\paragraph{Hashed retrieval.} Fix an order $n \in \{2,\dots,n_{\max}\}$ and form the suffix
$n$-gram $g_{t,n} = (x'_{t-n+1},\dots,x'_t)$. For each of $K$ hash heads $j = 1,\dots,K$ compute an index
into a table $\bm{E}_{n,j}$ with a prime number $P_{n,j}$ of rows, using the multiplicative--XOR hash
of \cref{eq:c10-mixhash}, and read one row:
\begin{equation}
  z_{t,n,j} = \varphi_{n,j}(g_{t,n}) \in \{0,\dots,P_{n,j}-1\},\qquad
  \ve_{t,n,j} = \bm{E}_{n,j}[z_{t,n,j}] \in \R^{d_h} .
  \label{eq:c10-engram-lookup}
\end{equation}
The memory vector is the concatenation of all rows,
\begin{equation}
  \ve_t = \big\Vert_{n=2}^{n_{\max}} \big\Vert_{j=1}^{K} \ve_{t,n,j} \in \R^{d_{\mathrm{mem}}},
  \qquad d_{\mathrm{mem}} = (n_{\max}-1) K d_h .
  \label{eq:c10-engram-concat}
\end{equation}
Positions $t < n$ are padded with a pad id. Nothing here depends on $\vh_t$: the indices are a
deterministic function of the token ids.

\begin{example}[A tiny hash]
Take multipliers $m_0 = 3$, $m_1 = 5$ and tables of $7$ and $11$ rows. The bigram $(x_{t-1}, x_t) =
(2, 4)$ mixes to $(4 \cdot 3) \oplus (2 \cdot 5) = 12 \oplus 10 = 1100_2 \oplus 1010_2 = 0110_2 = 6$,
so the two heads read rows $6 \bmod 7 = 6$ and $6 \bmod 11 = 6$. The reversed bigram $(4, 2)$ mixes
to $(2\cdot 3) \oplus (4 \cdot 5) = 00110_2 \oplus 10100_2 = 10010_2 = 18$ and reads rows $4$ and
$7$: order matters.
\end{example}

\paragraph{Context-aware gating.} The retrieved $\ve_t$ is a context-independent prior. It may be
wrong for the current context, because of a hash collision or because the $n$-gram is ambiguous.
Engram therefore uses the hidden state, which has already gathered context through the preceding
attention layers, as a query, and the memory as key and value:
\begin{equation}
  \vk_t = \mW_K \ve_t,\quad \vv_t = \mW_V \ve_t,\qquad
  g_t = \sigmoid\!\left(\frac{\RMSNorm(\vh_t)\T \RMSNorm(\vk_t)}{\sqrt{d}}\right),\qquad
  \tilde{\vv}_t = g_t\, \vv_t ,
  \label{eq:c10-engram-gate}
\end{equation}
with $\mW_K, \mW_V \in \R^{d \times d_{\mathrm{mem}}}$. (The paper writes $\alpha_t$ for the gate; we
write $g_t$ because $\alpha_t$ denotes a decay gate in this book.) This is attention with a single
candidate: instead of a softmax over many keys, a sigmoid decides between using the one retrieved
value and ignoring it. If the memory contradicts the context, $g_t \to 0$ suppresses it.

Why normalize both sides? Ignoring the learned scale and the small $\epsilon$,
\begin{gather*}
  \RMSNorm(\va) = \frac{\va}{\sqrt{\tfrac1d\sum_i a_i^2}}
  \quad\Longrightarrow\quad \norm{\RMSNorm(\va)}_2^2 = \frac{\norm{\va}_2^2}{\tfrac1d\norm{\va}_2^2} = d,\\
  \abs{\RMSNorm(\vh)\T\RMSNorm(\vk)} \le \norm{\RMSNorm(\vh)}_2\,\norm{\RMSNorm(\vk)}_2 = d
\end{gather*}
by the Cauchy--Schwarz inequality, so the gate logit lies in $[-\sqrt{d}, \sqrt{d}]$ whatever the
scale of $\vh_t$ and $\vk_t$. Without normalization a
large hidden state would saturate the sigmoid and stop the gradient; the paper cites gradient
stability as the reason.

\paragraph{Short convolution and residual.} Stack the gated values into
$\tilde{\mV} \in \R^{T \times d}$. A depthwise causal convolution with kernel size $w = 4$ and
dilation $\delta = n_{\max}$ widens the receptive field and adds nonlinearity:
\begin{equation}
  \mY = \SiLU\bigl(\mathrm{Conv1D}(\RMSNorm(\tilde{\mV}))\bigr) + \tilde{\mV},\qquad
  \mH \gets \mH + \mY ,
  \label{eq:c10-engram-conv}
\end{equation}
after which the block's attention and MoE run as usual. Row $t$ of the depthwise convolution is
$\sum_{i=0}^{w-1} \vw_i \had \vu_{t - i\delta}$ (channel-wise weights $\vw_i$, zero for negative
indices), so it sees positions $t, t-\delta, \dots, t-(w-1)\delta$. The paper initializes the
convolution to zero, so that at the start of training $\mY = \tilde{\mV}$.

\paragraph{Per-token and parallel views.} Everything in \crefrange{eq:c10-engram-lookup}{eq:c10-engram-gate}
is position-wise, and the convolution is a fixed causal filter, so the parallel computation over a
whole sequence and a token-by-token computation coincide; during decoding the module keeps only
the last $n_{\max}-1$ token ids and the last $(w-1)\delta$ normalized gated values per branch, a
state of constant size.

\begin{proposition}[Receptive field of an Engram layer]\label{prop:c10-engram-rf}
Hold the hidden states $\mH$ fixed. Then $\mY_t$ depends on the tokens
$x_{t-(n_{\max}-1)-(w-1)\delta}, \dots, x_t$ only. In particular, changing token $x_s$ can change
$\mY_t$ only for $s \le t \le s + (n_{\max}-1) + (w-1)\delta$.
\end{proposition}
\begin{proof}
Token $x_s$ enters the $n$-grams $g_{t,n}$ with $t - n + 1 \le s \le t$, i.e.\ positions
$t \in [s, s + n_{\max} - 1]$; these determine $\ve_t$, hence $\tilde{\vv}_t$. The convolution
output at $t$ reads $\tilde{\vv}_{t - i\delta}$ for $i = 0,\dots,w-1$, which adds a lag of up to
$(w-1)\delta$. The residual term only uses $\tilde{\vv}_t$.
\end{proof}
With $n_{\max} = 3$, $w = 4$, $\delta = 3$ this is a window of $1 + 2 + 9 = 12$ tokens. (The full
model of course also sees the whole context through $\vh_t$ in the gate.)

\paragraph{Multi-branch integration.} In the paper, the backbone uses manifold-constrained
hyper-connections, mHC for short \citep{c10-xie2025mhc,c10-zhu2025hyperconnections}. They widen the
residual stream into $R$ parallel branches ($R = 4$ in all experiments). Engram
shares one table and one $\mW_V$ across branches but gives each branch its own key projection, so
each branch has its own gate:
\begin{equation}
  g_t^{(r)} = \sigmoid\!\left(\frac{\RMSNorm(\vh^{(r)}_t)\T \RMSNorm(\mW^{(r)}_K \ve_t)}{\sqrt{d}}\right),
  \qquad \vu^{(r)}_t = g_t^{(r)}\, \mW_V \ve_t, \qquad r = 1,\dots,R .
  \label{eq:c10-engram-mhc}
\end{equation}
The $R+1$ projections can be fused into one dense matrix multiplication of shape
$(R+1)d \times d_{\mathrm{mem}}$, which the paper runs in FP8.

\begin{figure}[t]
\centering
\begin{tikzpicture}[font=\small, node distance=4mm,
  box/.style={draw,rounded corners=2pt,inner sep=3pt,align=center,minimum height=6mm},
  mem/.style={box,fill=orange!15}, op/.style={box,fill=blue!6}]
  \node[box] (ids) {token ids $x_{t-2}, x_{t-1}, x_t$};
  \node[op, right=6mm of ids] (cmp) {compress\\ $x' = \mathcal{P}(x)$};
  \node[op, right=6mm of cmp] (hash) {hash heads\\ $z_{t,n,j} = \varphi_{n,j}(g_{t,n})$};
  \node[mem, below=6mm of hash] (tab) {tables $\bm{E}_{n,j}$\\ (host or HBM)};
  \node[op, left=6mm of tab] (cat) {concat\\ $\ve_t \in \R^{d_{\mathrm{mem}}}$};
  \node[op, below=6mm of cat] (kv) {$\vk_t = \mW_K\ve_t$, $\vv_t = \mW_V\ve_t$};
  \node[box, left=6mm of kv] (h) {hidden $\vh_t$};
  \node[op, below=6mm of kv] (gate) {$g_t = \sigmoid(\RMSNorm(\vh_t)\T\RMSNorm(\vk_t)/\sqrt{d})$};
  \node[op, below=6mm of gate] (conv) {$\mY = \SiLU(\mathrm{Conv}(\RMSNorm(g\,\vv))) + g\,\vv$};
  \node[box, right=6mm of conv] (res) {$\mH \gets \mH + \mY$};
  \draw[-{Stealth}] (ids) -- (cmp); \draw[-{Stealth}] (cmp) -- (hash);
  \draw[-{Stealth}] (hash) -- (tab); \draw[-{Stealth}] (tab) -- (cat);
  \draw[-{Stealth}] (cat) -- (kv); \draw[-{Stealth}] (kv) -- (gate);
  \draw[-{Stealth}] (h.south) |- ([yshift=0mm]gate.west); \draw[-{Stealth}] (gate) -- (conv);
  \draw[-{Stealth}] (conv) -- (res);
  \node[font=\footnotesize, right=2mm of tab, align=left] {indices known\\ before the forward pass};
\end{tikzpicture}
\caption{Data flow of one Engram layer (single residual branch). The retrieval phase (top) depends
only on token ids; the fusion phase (bottom) uses the hidden state as the query of a scalar gate.}
\label{fig:c10-engram}
\end{figure}

\subsection{A new axis of sparsity: allocating parameters between MoE and memory}

The paper measures model size with three quantities: $P_{\mathrm{tot}}$, all trainable parameters
except the vocabulary embedding and output head; $P_{\mathrm{act}}$, the parameters activated per
token, which determines the training FLOPs; and the inactive budget
$P_{\mathrm{sparse}} = P_{\mathrm{tot}} - P_{\mathrm{act}}$ --- parameters that add capacity but no
compute, such as unselected experts or unretrieved table rows. An \emph{allocation ratio}
$\rho \in [0,1]$ splits the inactive budget:
\begin{equation}
  P^{(\mathrm{sparse})}_{\mathrm{MoE}} = \rho\, P_{\mathrm{sparse}},\qquad
  P_{\mathrm{Engram}} = (1-\rho)\, P_{\mathrm{sparse}} .
  \label{eq:c10-rho}
\end{equation}
$\rho = 1$ is a pure MoE model; $\rho < 1$ removes routed experts and spends the freed parameters on
table rows, keeping $P_{\mathrm{tot}}$ and $P_{\mathrm{act}}$ fixed. Engram rows are ``free'' in
$P_{\mathrm{act}}$ because each token reads a constant number of rows regardless of table size; the
small projections $\mW_K, \mW_V$ are counted as active.

Two regimes are studied. Under a fixed budget, validation loss as a function of $\rho$ is
\emph{U-shaped}: with $\rho \to 1$ the model has no dedicated store for static patterns and must
rebuild them with computation; with $\rho \to 0$ it lacks conditional computation for
context-dependent reasoning, and memory cannot replace it. The optimum lies at an intermediate
$\rho$ and was stable across the two compute scales tried. In the ``infinite memory'' regime, a
fixed MoE backbone receives ever larger tables, and the loss falls approximately linearly in the
logarithm of the number of slots --- a predictable scaling knob that costs no extra FLOPs. The
numbers are in the results box.

\paragraph{Per-token cost (our arithmetic).} One Engram layer costs, per token, $(n_{\max}-1)K$ row
reads, the fused projections ($2(R+1) d\, d_{\mathrm{mem}}$ floating-point operations), and
$\bigO(R d w)$ for the convolution and norms. With the paper's Engram-27B settings
($d = 2560$, $d_{\mathrm{mem}} = 1280$, $R = 4$) the projections are about $3.3 \times 10^7$ FLOPs
per layer, so the two Engram layers add well under one percent to the roughly
$2 \times 3.8 \times 10^9$ FLOPs per token of a model with $3.8$B activated parameters --- while the
tables hold billions of parameters.

\subsection{System design: offloading and prefetching}

Because the indices depend only on token ids, Engram's memory traffic is predictable, unlike MoE
routing, which needs the hidden state of the current layer. The paper exploits this in both phases
of the model's life.
\begin{itemize}
  \item \textbf{Training.} The tables are sharded across GPUs like any model-parallel parameter. An
        All-to-All exchange gathers the active rows in the forward pass and sends their gradients
        back in the backward pass, so total table capacity grows linearly with the number of
        accelerators. The embedding rows are optimized with Adam at a $5\times$ learning rate and no
        weight decay, while the backbone uses Muon.
  \item \textbf{Inference.} The tables live in host DRAM. As soon as the token ids are known, the
        host gathers the needed rows and transfers them over PCIe \emph{asynchronously}, overlapping
        the transfer with the computation of the Transformer blocks that precede the Engram layer.
        Communication volume scales with the number of activated rows, not with the table size.
  \item \textbf{Placement.} Putting Engram deeper gives more compute time to hide the transfer, but
        modeling quality prefers early insertion, where static patterns can be offloaded before the
        backbone spends depth on them. The paper's models use layers 2 and 15 of a 30-block network:
        an early layer for modeling, a later one that also benefits from richer gating queries.
  \item \textbf{Caching.} $N$-gram frequencies follow a Zipf distribution, so a small fraction of
        rows receives most accesses. Frequent rows can be cached in HBM or DRAM and the long tail
        kept on slower storage such as NVMe, a multi-level hierarchy proposed in the paper.
\end{itemize}

\subsection{Algorithms and complexity}

\begin{algorithm}[t]
\caption{Engram address computation (host side, before or during the forward pass)}\label{alg:c10-engram-hash}
\begin{algorithmic}[1]
\Require token ids $x_{1:T}$, compression table $\mathcal{P}$, multipliers $m^{(\ell)}_{0:n_{\max}-1}$ and primes $P^{(\ell)}_{n,j}$ for each Engram layer $\ell$
\Ensure row indices $z^{(\ell)}_{t,n,j}$
\State $x'_t \gets \mathcal{P}(x_t)$ for all $t$; pad positions $t \le 0$ with the compressed pad id
\For{each Engram layer $\ell$, each order $n = 2,\dots,n_{\max}$}
  \State $\mathrm{mix}_t \gets \bigoplus_{i=0}^{n-1} x'_{t-i}\, m^{(\ell)}_i$ for all $t$ \Comment{vectorized over $t$}
  \For{$j = 1,\dots,K$}
    \State $z^{(\ell)}_{t,n,j} \gets \mathrm{mix}_t \bmod P^{(\ell)}_{n,j}$; issue prefetch of $\bm{E}^{(\ell)}_{n,j}[z^{(\ell)}_{t,n,j}]$
  \EndFor
\EndFor
\end{algorithmic}
\end{algorithm}

\begin{algorithm}[t]
\caption{Engram layer forward (one residual branch; repeat the gate per branch for mHC)}\label{alg:c10-engram-layer}
\begin{algorithmic}[1]
\Require hidden states $\mH \in \R^{T \times d}$, fetched rows $\ve_{t,n,j}$
\State $\ve_t \gets \Vert_{n,j}\, \ve_{t,n,j}$ \Comment{$\R^{d_{\mathrm{mem}}}$}
\State $\vk_t \gets \mW_K\ve_t$;\; $\vv_t \gets \mW_V\ve_t$
\State $g_t \gets \sigmoid\bigl(\RMSNorm(\vh_t)\T\RMSNorm(\vk_t)/\sqrt{d}\bigr)$;\; $\tilde{\vv}_t \gets g_t \vv_t$
\State $\mY \gets \SiLU(\mathrm{CausalDepthwiseConv}_{w,\delta}(\RMSNorm(\tilde{\mV}))) + \tilde{\mV}$
\State \Return $\mH + \mY$ \Comment{then attention and MoE of the same block}
\end{algorithmic}
\end{algorithm}

\paragraph{Complexity.} Address computation is $\bigO(T n_{\max}^2 K)$ integer operations (or
$\bigO(T n_{\max} K)$ reusing partial mixes), independent of the table sizes. The forward pass
costs $\bigO(T (R+1) d\, d_{\mathrm{mem}})$ for the projections, $\bigO(T R d w)$ for the
convolution and $\bigO(T d_{\mathrm{mem}})$ of memory traffic. The parameters of one Engram layer
are $\sum_{n,j} P_{n,j}\, d_h$ table entries plus $(R+1) d\, d_{\mathrm{mem}}$ projection weights.
In decoding, the per-token cost and the state are constant in the sequence length.

\subsection{Implementation walkthrough}\label{sec:c10-engram-code}

The official repository \repo{deepseek-ai__Engram} contains the paper
(\repofile{repos/deepseek-ai__Engram/Engram_paper.pdf}), figures, a draw.io architecture diagram and
a single script, \repofile{repos/deepseek-ai__Engram/engram_demo_v1.py}. Its README warns that the
script is ``a demonstration version intended to illustrate the data flow'' that mocks attention,
MoE and the hyper-connections. It is still the best place to see the exact hashing and fusion.

\begin{implnote}[title={In the code: configuration}]
Two dataclasses fix the sizes. \code{EngramConfig}: tokenizer \code{deepseek-ai/DeepSeek-V3};
\code{engram_vocab_size = [129280*5, 129280*5]} (the requested table size per head for 2-grams and
3-grams); \code{max_ngram_size = 3}; \code{n_embed_per_ngram = 512}; \code{n_head_per_ngram = 8};
\code{layer_ids = [1, 15]}; \code{pad_id = 2}; \code{seed = 0}; \code{kernel_size = 4}.
\code{BackBoneConfig}: \code{hidden_size = 1024}, \code{hc_mult = 4} (the $R = 4$ branches),
\code{vocab_size = 129280}, \code{num_layers = 30}. So $d_h = 512/8 = 64$ and
$d_{\mathrm{mem}} = 2 \cdot 512 = 1024$. In our arithmetic, the 16 prime-sized tables of one Engram
layer have about $1.03 \times 10^7$ rows, i.e.\ about $6.6 \times 10^8$ parameters per layer at
$d_h = 64$.
\end{implnote}

\begin{implnote}[title={In the code: tokenizer compression and hash parameters}]
\code{CompressedTokenizer} decodes every token id of the base tokenizer to text, normalizes it
with a Hugging Face \code{normalizers.Sequence} (NFKC, NFD, StripAccents, Lowercase, collapse of
whitespace runs to one space, Strip, with a sentinel so that a pure-space token survives), and
assigns one new id per distinct normalized string; tokens that decode to the Unicode replacement
character keep their raw token string as key. The result is a NumPy array \code{lookup_table}
mapping old to new ids. \code{NgramHashMapping} draws, for each Engram layer, multipliers
\code{r * 2 + 1} (odd) from a generator seeded with \code{seed + 10007 * layer_id}, with
\code{r} below \code{half_bound = (int64_max // V') // 2}. Then every product of a compressed id
($< V'$) with a multiplier ($< 2\cdot$\code{half_bound}) stays below the \code{int64} maximum,
so the mix never overflows. \code{calculate_vocab_size_across_layers} searches upwards from the
requested size for primes, keeping a global \code{seen_primes} set so that \emph{every} head of
every order in every layer gets a distinct prime size --- the condition of \cref{prop:c10-crt}.
\end{implnote}

\begin{lstlisting}[caption={Excerpt from \texttt{repos/deepseek-ai\_\_Engram/engram\_demo\_v1.py} (\texttt{NgramHashMapping.\_get\_ngram\_hashes}).},label={lst:c10-engram-hash}]
        x = np.asarray(input_ids, dtype=np.int64)
        B, T = x.shape

        multipliers = self.layer_multipliers[layer_id]

        def shift_k(k: int) -> np.ndarray:
            if k == 0: return x
            shifted = np.pad(x, ((0, 0), (k, 0)),
                                mode='constant', constant_values=self.pad_id)[:, :T]
            return shifted

        base_shifts = [shift_k(k) for k in range(self.max_ngram_size)]

        all_hashes = []

        for n in range(2, self.max_ngram_size + 1):
            n_gram_index = n - 2
            tokens = base_shifts[:n]
            mix = (tokens[0] * multipliers[0])
            for k in range(1, n):
                mix = np.bitwise_xor(mix, tokens[k] * multipliers[k])
            num_heads_for_this_ngram = self.n_head_per_ngram
            head_vocab_sizes = self.vocab_size_across_layers[layer_id][n_gram_index]

            for j in range(num_heads_for_this_ngram):
                mod = int(head_vocab_sizes[j])
                head_hash = mix % mod
                all_hashes.append(head_hash.astype(np.int64, copy=False))

        return np.stack(all_hashes, axis=2)
\end{lstlisting}

The shifted copies \code{base_shifts[k]} hold $x'_{t-k}$ (left-padded with the pad id), so
\code{mix} is exactly \cref{eq:c10-mixhash}, computed for all positions at once; the result has
shape $(B, T, (n_{\max}-1)K) = (B, T, 16)$. \code{MultiHeadEmbedding} stores all 16 tables of a
layer in a single \code{nn.Embedding} with $\sum_{n,j} P_{n,j}$ rows and adds per-head offsets to the
indices before the lookup, a standard trick that turns 16 gathers into one.

\begin{lstlisting}[caption={Excerpt from \texttt{repos/deepseek-ai\_\_Engram/engram\_demo\_v1.py} (\texttt{Engram.forward}).},label={lst:c10-engram-fwd}]
    def forward(self,hidden_states,input_ids):
        """
        hidden_states: [B, L, HC_MULT, D]
        input_ids: [B, L]
        """
        hash_input_ids = torch.from_numpy(self.hash_mapping.hash(input_ids)[self.layer_id])
        embeddings = self.multi_head_embedding(hash_input_ids).flatten(start_dim=-2)
        gates = []
        for hc_idx in range(backbone_config.hc_mult):
            key = self.key_projs[hc_idx](embeddings)
            normed_key = self.norm1[hc_idx](key)
            query = hidden_states[:,:,hc_idx,:]
            normed_query = self.norm2[hc_idx](query)
            gate = (normed_key * normed_query).sum(dim=-1) / math.sqrt(backbone_config.hidden_size)
            gate = gate.abs().clamp_min(1e-6).sqrt() * gate.sign()
            gate = gate.sigmoid().unsqueeze(-1)
            gates.append(gate)
        gates = torch.stack(gates,dim=2)
        value = gates * self.value_proj(embeddings).unsqueeze(2)
        output = value + self.short_conv(value)
        return output
\end{lstlisting}

\begin{table}[t]
\centering
\small
\caption{Tensor shapes in \texttt{Engram.forward} of the demo ($B$ batch, $T$ length,
$R = 4$ branches, $d = 1024$, $d_{\mathrm{mem}} = 1024$).}
\label{tab:c10-engram-shapes}
\begin{tabular}{@{}lll@{}}
\toprule
Tensor & Shape & Meaning \\
\midrule
\code{input_ids} & $(B, T)$ & raw token ids (hashed on the CPU with NumPy) \\
\code{hash_input_ids} & $(B, T, 16)$ & one row index per (order, head) \\
\code{embeddings} & $(B, T, 16, 64) \to (B, T, 1024)$ & $\ve_t$ after flattening \\
\code{key} (per branch) & $(B, T, 1024)$ & $\mW^{(r)}_K \ve_t$ \\
\code{hidden_states} & $(B, T, 4, 1024)$ & $\vh^{(r)}_t$ for the 4 branches \\
\code{gates} & $(B, T, 4, 1)$ & $g^{(r)}_t$ \\
\code{value} & $(B, T, 4, 1024)$ & $g^{(r)}_t \mW_V \ve_t$ (shared $\mW_V$ broadcast) \\
\code{output} & $(B, T, 4, 1024)$ & $\mY$, added to every branch by the caller \\
\bottomrule
\end{tabular}
\end{table}

\begin{implnote}[title={In the code: differences between the demo and the paper}]
(i) The gate has an extra step: after the scaled dot product, line
\code{gate.abs().clamp_min(1e-6).sqrt() * gate.sign()} takes a signed square root before the
sigmoid. This does not appear in the paper's equation. Since the logit lies in
$[-\sqrt{d}, \sqrt{d}]$, the signed square root compresses it to $[-d^{1/4}, d^{1/4}]$ ($\pm 5.7$
for $d = 1024$), which keeps the sigmoid away from saturation; this reading is ours.
(ii) \code{ShortConv} applies a per-branch \code{nn.RMSNorm}, concatenates the branches into
$R d$ channels, runs \code{nn.Conv1d} with \code{groups} equal to the channel count (depthwise),
\code{padding=(kernel_size - 1) * dilation} and truncation to the first $T$ outputs (which makes it
causal), and SiLU; the dilation is \code{max_ngram_size} $= 3$, as in the paper. The demo does not
zero-initialize the convolution, which the paper does.
(iii) \code{TransformerBlock} applies Engram as a residual \emph{before} attention and MoE, which
are identity lambdas here; the hyper-connections are mocked by copying the embedding into four
branches and reading branch 0 at the end. (iv) The hashes are recomputed with NumPy on the CPU
inside every Engram layer's forward pass; a production system would compute them once, ahead of
time, and prefetch the rows (\cref{alg:c10-engram-hash}). (v) The demo's \code{layer_ids = [1, 15]}
are 0-based indices; the paper's Engram-27B uses layers 2 and 15.
\end{implnote}

\begin{lstlisting}[caption={Reference implementation (ours): hashed $N$-gram lookup with context-aware gated fusion (one branch).},label={lst:c10-engram-ref}]
import torch, torch.nn as nn, torch.nn.functional as F

class EngramLayer(nn.Module):
    def __init__(self, d, primes, mults, d_head, pad_id, w=4):
        """primes[n][j]: rows of table (n, j); mults: odd ints, one per n-gram position."""
        super().__init__()
        self.primes, self.mults, self.pad = primes, mults, pad_id
        sizes = [p for n in sorted(primes) for p in primes[n]]
        self.register_buffer("offsets", torch.tensor([0] + sizes[:-1]).cumsum(0))
        self.table = nn.Embedding(sum(sizes), d_head)          # all heads in one table
        d_mem = len(sizes) * d_head
        self.w_k, self.w_v = nn.Linear(d_mem, d, bias=False), nn.Linear(d_mem, d, bias=False)
        self.n_h, self.n_k, self.n_y = nn.RMSNorm(d), nn.RMSNorm(d), nn.RMSNorm(d)
        self.delta = max(primes)                                # dilation = max order
        self.conv = nn.Conv1d(d, d, w, groups=d, bias=False, dilation=self.delta,
                              padding=(w - 1) * self.delta)
        nn.init.zeros_(self.conv.weight)                        # Y = gated value at init

    def indices(self, ids):                                     # ids: (B, T) compressed
        T, cols = ids.shape[1], []
        shift = lambda k: F.pad(ids, (k, 0), value=self.pad)[:, :T]
        for n in sorted(self.primes):
            mix = shift(0) * self.mults[0]
            for k in range(1, n):
                mix = torch.bitwise_xor(mix, shift(k) * self.mults[k])
            cols += [mix % p for p in self.primes[n]]
        return torch.stack(cols, -1) + self.offsets             # (B, T, heads)

    def forward(self, h, ids):                                  # h: (B, T, d)
        e = self.table(self.indices(ids)).flatten(-2)           # (B, T, d_mem)
        k, v = self.w_k(e), self.w_v(e)
        g = torch.sigmoid((self.n_h(h) * self.n_k(k)).sum(-1, keepdim=True) / h.shape[-1] ** 0.5)
        vt = g * v                                              # gated value
        u = self.conv(self.n_y(vt).transpose(1, 2))[..., : h.shape[1]].transpose(1, 2)
        return h + F.silu(u) + vt                               # residual update
\end{lstlisting}

\paragraph{Numerical checks.} In NumPy (scratch space) we checked that (a) a vectorized
multiplicative--XOR hash over shifted id arrays, as in \cref{lst:c10-engram-hash}, equals a naive
loop over positions and yields only non-negative indices; (b) the all-head collision property of
\cref{prop:c10-crt}, with the demo-sized primes; (c) the receptive field of
\cref{prop:c10-engram-rf}: with $n_{\max} = 3$, $w = 4$, $\delta = 3$, changing token 20 of a
40-token sequence changed the outputs at positions 20 to 31 and nowhere else; and (d) parameter
arithmetic. Reading Table~5 of the paper, ``Engram Vocab Size'' $2{,}262{,}400$ (Engram-27B) and
$7{,}239{,}680$ (Engram-40B) with $d_{\mathrm{mem}} = 1280$, 16 heads of width 80 and two Engram
layers, the interpretation ``rows per head'' gives $2 \times 2{,}262{,}400 \times 1280 \approx
5.79 \times 10^9$ and $18.53 \times 10^9$ table parameters, close to the reported 5.7B and 18.5B.
This interpretation is ours; the paper does not spell out the formula.

\begin{results}
\textbf{Source file:} \repofile{repos/deepseek-ai__Engram/Engram_paper.pdf} (the README shows the
same figures as images).
\begin{itemize}
  \item \textbf{Iso-parameter, iso-FLOPs pre-training} (Table~1; 262B tokens, 3.8B activated
        parameters, MoE-27B with 72 routed experts versus Engram-27B with 55 routed experts and a
        5.7B-parameter Engram): validation loss $1.634 \to 1.622$; MMLU $57.4 \to 60.4$; CMMLU
        $57.9 \to 61.9$; BBH $50.9 \to 55.9$; ARC-Challenge $70.1 \to 73.8$; DROP $55.7 \to 59.0$;
        HumanEval $37.8 \to 40.8$; GSM8K $58.4 \to 60.6$; MATH $28.3 \to 30.7$. Engram-40B (18.5B
        Engram parameters, same activated budget) reaches validation loss $1.610$.
  \item \textbf{Allocation} (Section~3.1): at $6\times10^{20}$ FLOPs, validation loss improves from
        $1.7248$ at $\rho = 100\%$ to $1.7109$ near $\rho \approx 80\%$; the optimum stays at
        $\rho \approx 75$--$80\%$ across the two compute regimes.
  \item \textbf{Ablation} (Section~6.2; 12-layer 3B MoE, 100B tokens): baseline validation loss
        $1.808$; with a 1.6B-parameter Engram at layers 2 and 6, $1.768$; the best single
        insertion layer is layer 2 ($1.770$).
  \item \textbf{Long context} (Table~2, iso-loss comparison of Engram-27B at 46k steps with MoE-27B
        at 50k): RULER multi-query needle-in-a-haystack $84.2 \to 97.0$, variable tracking
        $77.0 \to 87.2$.
  \item \textbf{Offloading} (Table~4, NVIDIA H800, a 100B-parameter table in host memory):
        throughput 9{,}031.62 $\to$ 8{,}858.28 tokens/s on a 4B dense backbone and 6{,}315.52 $\to$
        6{,}140.02 tokens/s on an 8B dense backbone.
  \item \textbf{Sensitivity} (Section~6.3): with the Engram output suppressed at inference,
        factual-knowledge benchmarks retain only 29--44\% of their scores (TriviaQA 29\%) while
        reading comprehension retains 81--93\% (C3 93\%).
\end{itemize}
\end{results}

\subsection{Discussion}

Engram makes three points that go beyond $N$-gram embeddings. First, a lookup memory is not a
replacement for computation but a complement to it: the allocation curve is U-shaped, so neither
pure MoE nor a memory-heavy model is best. Second, the benefit is not confined to knowledge-heavy
tasks: the paper's LogitLens and CKA analyses suggest that Engram relieves early layers from
reconstructing static patterns, which acts like extra effective depth, and that delegating local
dependencies frees attention for long-range context. Third, deterministic addressing turns a
modeling choice into a systems advantage: tables can be sharded in training and offloaded with
prefetching in inference, which is impossible for hidden-state-dependent routing.

The limitations follow from the same design. The memory is static at inference: there is no write
path, and editing a fact would mean editing hashed rows shared by colliding $n$-grams. The key is a
short window of at most $n_{\max}$ tokens, so anything that needs longer context must come from
$\vh_t$ through the gate. Collisions are absorbed by multi-head hashing and gating rather than
avoided. The public code is a demonstration, not the training system. Relative to the other systems
here, Engram is the hashed, sparse cousin of the embedding table, and MemoryLLM
(\cref{sec:c10-memoryllm}) is its unigram, every-layer cousin. Hybrid models that combine several
memory types are discussed in \cref{ch:hybrid}; the survey mentions Engram-inspired lookup
augmentations among them.

% =====================================================================================
\section{ExplicitLM: a bank of human-readable memories}\label{sec:c10-explicitlm}

\subsection{Motivation}

Knowledge stored in weights is entangled and not addressable: one cannot point at the parameters
that hold a fact, check them, or replace them when the fact changes. ExplicitLM
\citep{yu2025explicitlm} takes the editable-memory direction of PlugLM (\cref{sec:c10-precursors})
to its literal end. According to its abstract, the model has a million-scale memory bank in which
\emph{each entry is a human-readable token sequence}, so the knowledge base can be inspected and
modified directly. To make such a bank usable inside a network it needs (i) an address step that is
cheap for a million entries --- product keys --- and (ii) a selection step that is discrete (one
entry, readable) yet trainable end to end --- the Gumbel-Softmax with a straight-through
estimator. The abstract also describes a knowledge partition inspired by dual-system cognitive
theory: a frozen part holding explicit facts and a learnable part holding implicit patterns
(the abstract gives a 20\%/80\% split), maintained with exponential-moving-average (EMA) updates.

The survey stresses that such systems are not retrieval-augmented generation: the memory banks are
trained and invoked as architectural components rather than appended through an external
orchestration pipeline, and their significance is that they make knowledge more inspectable and
expose targeted update operations (\S IV-B(d), (f)).

\begin{taxonomy}
\textbf{Representation:} explicit --- layer-wise memory banks whose entries are token sequences
\citep[\S IV-B(d)]{zhoubian2026memory}; Table~I's paradigm label is ``Interpretable Memory Banks''.
\textbf{Update dynamics:} offline (Table~I). The repository version additionally rewrites entries at
inference time (see the implementation notes).
\textbf{Persistence:} long-term.
\textbf{Update rule:} not listed in Table~II. In our reading the paper's EMA maintenance of
learnable entries is objective-induced, and the repository's inference-time rewriting is an
admission/eviction rule (FIFO, LRU and similar policies).
\end{taxonomy}

\subsection{Formulation}

The abstract fixes the ingredients; the repository fixes many details. Where they may differ, we
follow the code and say so.

\paragraph{The bank.} The memory holds $N = n^2$ entries. Entry $i$ is a sequence of $\ell_m$ token
ids, $m_i = (m_{i,1},\dots,m_{i,\ell_m}) \in \mathcal{V}^{\ell_m}$, padded or truncated to a fixed
length ($N = 2^{20}$ and $\ell_m = 16$ by default in the code). To use an entry inside the network
it needs a vector. The code embeds each token with the model's own embedding table and averages:
\begin{equation}
  \bar{\ve}_i = \frac{1}{\ell_m}\sum_{j=1}^{\ell_m} \mathrm{Emb}(m_{i,j}) \in \R^{d} .
  \label{eq:c10-entry}
\end{equation}
So the memory's contents are \emph{not} free parameters: they are fully determined by token ids,
which is exactly what makes them readable (decode the ids) and editable (overwrite the ids).

\paragraph{Stage 1: coarse filtering with product keys.} A query network $f_{\vtheta}$ maps a hidden
state to $\vq \in \R^{d_k}$. The two halves of the query are compared with two sets of
$n = \sqrt{N}$ sub-keys by cosine similarity with a learned temperature $\tau$:
\begin{equation}
  a_i = \tau\, \cos\bigl(\vq^{(1)}, \bm{c}^{(1)}_i\bigr),\qquad
  b_j = \tau\, \cos\bigl(\vq^{(2)}, \bm{c}^{(2)}_j\bigr),\qquad s_{ij} = a_i + b_j ,
  \label{eq:c10-explicit-coarse}
\end{equation}
and the top-$c$ pairs over the $n\times n$ grid form the candidate set $\mathcal{C}$, found exactly
from shortlists of length $k' \ge c$ (\cref{prop:c10-pkm}). The abstract describes this, in its own notation, as reducing
the cost from $\bigO(N \cdot \abs{I})$ to $\bigO(\sqrt{N}\cdot\abs{I})$.

\paragraph{Stage 2: fine selection with Gumbel-Softmax.} For each token $t$ the model scores the
candidates by cosine similarity with its (normalized) hidden state,
$u_{t,i} = \cos(\vh_t, \bar{\ve}_i)$ for $i \in \mathcal{C}$, and selects one of them. A hard
$\argmax$ has no useful gradient, so ExplicitLM uses the Gumbel-Softmax relaxation
\citep{c10-jang2017gumbel,c10-maddison2017concrete}: draw independent $G_i \sim \mathrm{Gumbel}(0,1)$,
\begin{equation}
\begin{gathered}
  \vy^{\mathrm{soft}}_t = \softmax\bigl((\vu_t + \vg)/\tau_g\bigr),\qquad
  \vy^{\mathrm{hard}}_t = \mathrm{onehot}\bigl(\argmax_i (u_{t,i} + G_i)\bigr),\\
  \vy_t = \vy^{\mathrm{hard}}_t - \mathrm{sg}\bigl(\vy^{\mathrm{soft}}_t\bigr) + \vy^{\mathrm{soft}}_t ,
\end{gathered}
  \label{eq:c10-gumbel}
\end{equation}
where $\mathrm{sg}$ stops the gradient. In the forward pass $\vy_t$ equals the one-hot vector, so
exactly one readable entry is used; in the backward pass the gradient is that of the soft weights
--- the straight-through estimator \citep{c10-bengio2013ste}. The selected memory is
$\vm_t = \sum_{i\in\mathcal{C}} y_{t,i}\, \bar{\ve}_i$.

Why does adding Gumbel noise and taking the $\argmax$ produce a sample from the softmax? The
derivation is short and worth doing once.

\begin{proposition}[Gumbel-max trick]\label{prop:c10-gumbel}
Let $G_1,\dots,G_C$ be i.i.d.\ with CDF $F(g) = \exp(-e^{-g})$. Then
$\Pr\bigl[\argmax_i (u_i + G_i) = k\bigr] = e^{u_k}/\sum_i e^{u_i}$.
\end{proposition}
\begin{proof}
The density is $f(g) = e^{-g}\exp(-e^{-g})$. Index $k$ wins when $G_i < u_k + G_k - u_i$ for all
$i \ne k$. Conditioning on $G_k = g$ and using independence,
\[
  \Pr[k \text{ wins}] = \int_{-\infty}^{\infty} f(g) \prod_{i\ne k} F(u_k + g - u_i)\, dg
  = \int_{-\infty}^{\infty} e^{-g} \exp\Bigl(-e^{-g}\Bigl[1 + \sum_{i \ne k} e^{u_i - u_k}\Bigr]\Bigr) dg .
\]
The bracket equals $Z e^{-u_k}$ with $Z = \sum_i e^{u_i}$. Substitute $v = e^{-g}$, so
$dv = -e^{-g}dg$ and the limits $g = -\infty, \infty$ become $v = \infty, 0$:
$\Pr[k \text{ wins}] = \int_0^\infty \exp(-v Z e^{-u_k})\, dv = e^{u_k}/Z$.
\end{proof}

As $\tau_g \to 0$, $\vy^{\mathrm{soft}}_t$ approaches the one-hot sample, so the soft weights used
for the gradient become more faithful to the forward choice but also higher-variance. In our NumPy
check (scratch space), 400{,}000 Gumbel-max draws with logits $(1.0, 0.2, -0.5, 2.0, 0.0)$ matched
the softmax probabilities within $1.3$ standard errors in every coordinate, the straight-through
weights equalled the one-hot vector in the forward pass, and the mean largest relaxed weight rose
from $0.62$ at $\tau_g = 1$ to $0.99$ at $\tau_g = 0.03$.

\paragraph{Fusion.} The code fuses the selected memory with a SwiGLU-style gated MLP on the
concatenation $[\vh_t; \vm_t]$ and scales the result by a learned, similarity-dependent scalar:
\begin{equation}
  \vo_t = \mW_{\mathrm{down}}\Bigl(\SiLU\bigl(\mW_{\mathrm{gate}}[\vh_t;\vm_t]\bigr) \had
  \mW_{\mathrm{up}}[\vh_t;\vm_t]\Bigr),\qquad
  \vx_t \gets \vx_t + \lambda\, \vo_t ,
  \label{eq:c10-explicit-fuse}
\end{equation}
where $\lambda \in [0,1]$ is a clamped small MLP of the mean selected similarity plus a learned
bias (initialized to zero) --- one scalar per sequence. Because the update is additive, the
backbone still works if the memory contributes nothing (the code calls this a ``shortcut''
mechanism).

\paragraph{Auxiliary losses (code).} Besides the language-modeling loss, every memory layer
contributes a similarity loss $\loss_{\mathrm{sim}} = -\mathrm{mean}_t \sum_i y_{t,i} u_{t,i}$,
which rewards selecting candidates that match the hidden state, and a diversity loss
$\loss_{\mathrm{div}}$, the mean off-diagonal cosine similarity among the candidates' vectors,
which discourages redundant candidate sets. They are averaged over layers and weighted
$0.1$ and $0.05$ by default.

\paragraph{Updating the memory.} Two mechanisms are described, and they are different.
\begin{itemize}
  \item \textbf{Paper (abstract).} Entries are partitioned into frozen explicit facts and learnable
        implicit patterns, the latter maintained by EMA updates for stability. One way to formalize
        an EMA update of a learnable entry $i$ is
        $\tilde{\ve}_i \gets (1-\gamma)\tilde{\ve}_i + \gamma\, \bar{\vh}_i$, where $\bar{\vh}_i$ is the
        mean hidden state of the tokens that selected entry $i$ in the current batch and $\gamma$ a
        small rate; to keep the entry human-readable, the updated vector must then be mapped back to
        token ids (for example by nearest-neighbour decoding against the embedding table). The
        paper's exact parameterization may differ, and the mirrored code contains no EMA path.
  \item \textbf{Repository.} The bank is a buffer of token ids that is fixed during training and
        rewritten at inference: a fact extractor compresses incoming text with an LLMLingua-2
        model, splits it into sentences, tokenizes and pads them to $\ell_m$ tokens, and writes them
        into slots chosen by a replacement policy (FIFO, LRU, random or importance-based).
\end{itemize}

\subsection{Algorithms and complexity}

\begin{algorithm}[t]
\caption{ExplicitLM memory layer (two-stage retrieval and fusion)}\label{alg:c10-explicit}
\begin{algorithmic}[1]
\Require block output $\mX \in \R^{T\times d}$, bank $\{m_i\}_{i=1}^{N}$, sub-keys $\mC^{(1)},\mC^{(2)}$, shortlist $k'$, candidates $c$
\State $\mH \gets \RMSNorm(\mX)$
\State $\vq \gets f_{\vtheta}(\cdot)$ \Comment{per token (ours) or one per sequence (code: mean of $\mH$)}
\State $\va, \vb \gets$ scaled cosine scores of $\vq^{(1)}, \vq^{(2)}$ with the sub-keys
\State $\mathcal{C} \gets$ top-$c$ of $\{a_i + b_j\}$ over $\TopK(\va,k') \times \TopK(\vb,k')$ \Comment{\cref{prop:c10-pkm}}
\State $\bar{\ve}_i \gets \frac{1}{\ell_m}\sum_j \mathrm{Emb}(m_{i,j})$ for $i \in \mathcal{C}$ \Comment{gather $c\,\ell_m$ token rows}
\For{each token $t$}
  \State $u_{t,i} \gets \cos(\vh_t, \bar{\ve}_i)$ for $i \in \mathcal{C}$;\; $\vy_t \gets$ straight-through Gumbel-Softmax of $\vu_t$
  \State $\vm_t \gets \sum_i y_{t,i}\bar{\ve}_i$;\; $\vx_t \gets \vx_t + \lambda\,\mathrm{SwiGLU}([\vh_t;\vm_t])$
\EndFor
\State \Return $\mX$, auxiliary losses $\loss_{\mathrm{sim}}, \loss_{\mathrm{div}}$
\end{algorithmic}
\end{algorithm}

\paragraph{Cost.} Per query, stage 1 costs the query network plus $\bigO(\sqrt{N} d_k)$ for the
sub-key scores and $\bigO(k'^2)$ to merge the shortlists; stage 2 costs $\bigO(c\,\ell_m d)$ to build
the candidate vectors and $\bigO(T c\, d)$ for the similarities; fusion is $\bigO(T d^2)$. With the
code's defaults ($N = 2^{20}$, so $\sqrt{N} = 1024$, $k' = 128$, $c$ between 16 and 32) the coarse
stage scores 2048 sub-keys and merges $16{,}384$ pairs instead of scoring a million entries. The
bank itself is only token ids, $N \ell_m$ integers ($2^{24}$ of them by default), and the
``values'' share the model's embedding table, so a million-entry memory adds no value parameters
at all; the trainable parameters are the query networks, sub-keys (frozen in the code) and fusion
modules.

\subsection{Implementation walkthrough}\label{sec:c10-explicit-code}

\begin{implnote}[title={In the code: which version is this?}]
The mirror \repo{SCUT-HCC__ExplicitLM} comes from the GitHub organization of the senior author's
lab. Its README (in Chinese) never cites the paper, and it describes a system built on a
\emph{frozen Qwen3-4B} backbone in which only the memory components (\code{MemoryGate},
\code{GatedMemoryFusion}, \code{MemoryNorm}) are trained --- about 0.208B trainable parameters
according to the README. \repofile{repos/SCUT-HCC__ExplicitLM/CLEANUP_SUMMARY.md} records that an
older small-model version (width 512, 8 layers, a 6400-token vocabulary, a \code{1_pretrain.py}
pre-training script) was moved to an archive directory that is not in the mirror. The repository
may therefore describe a later version than the paper, and we treat it as ``probably official''.
Dataset files over 20\,MB were left out of the mirror.
\end{implnote}

The model is \code{ExplicitLM} in \repofile{repos/SCUT-HCC__ExplicitLM/models/core/ExplicitLM.py}.
It builds the Qwen3 embedding, rotary embedding and final norm, and replaces every decoder layer by
a \code{Qwen3ExplicitLMBlock} (\repofile{repos/SCUT-HCC__ExplicitLM/models/core/Qwen3ExplicitLMBlock.py}):
a standard \code{Qwen3DecoderLayer} followed by its own memory norm, memory gate and fusion. All 36
layers of Qwen3-4B (\repofile{repos/SCUT-HCC__ExplicitLM/config/qwen3_4b_params.py}) therefore
retrieve independently from the same bank. The bank is a persistent buffer of shape
$(N, \ell_m)$ holding token ids, initialized randomly and then filled from a knowledge base of
sentences (\code{sentence_trex_data.json}, i.e.\ T-REx-derived facts) by
\code{MemoryBankProcessor} in \repofile{repos/SCUT-HCC__ExplicitLM/utils/model_initializer.py}
(each sentence tokenized and padded or truncated to $\ell_m$ tokens, the result cached as a
\code{.pt} file). In the mirrored file, the call that builds the bank is indented inside the branch
that creates a missing cache directory, which looks like an indentation slip. The bank is kept on
the CPU and moved to the GPU when used.

\begin{implnote}[title={In the code: stage 1 (\texttt{MemoryGate})}]
\repofile{repos/SCUT-HCC__ExplicitLM/models/memory_bank/MemoryGate.py} implements the query network
as a linear layer, two pre-norm residual GELU blocks, a LayerNorm and a projection to
\code{knowledge_dim} $= 1536$. The sub-keys \code{keys} have shape $(2, \sqrt{N}, d_k/2)$; when
\code{keys_path} exists they are loaded and \emph{frozen}. Scores are cosine similarities times
\code{logit_scale.exp()}, a CLIP-style learnable temperature initialized to $1/0.07$ and clamped at
100. \code{generate_candidates} (\cref{lst:c10-explicit-cand}) is the product-key search of
\cref{lst:c10-pkm} with \code{num_candidates_internal} $= 128$ and \code{num_candidates} $= 32$ by
default (the training configuration uses 16). The LoRA variant
\repofile{repos/SCUT-HCC__ExplicitLM/models/memory_bank/MemoryGateLoRA.py} factorizes the linear
layers with rank \code{gate_rank} $= 128$ and keeps one residual block.
\end{implnote}

\begin{lstlisting}[caption={Excerpt from \texttt{repos/SCUT-HCC\_\_ExplicitLM/models/memory\_bank/MemoryGate.py} (\texttt{generate\_candidates}; comment lines in Chinese elided).},label={lst:c10-explicit-cand}]
        bsz, seq_len, _ = scores_1.shape
        # ...
        topk_scores_1, topk_indices_1 = scores_1.topk(self.num_candidates_internal, dim=-1)
        topk_scores_2, topk_indices_2 = scores_2.topk(self.num_candidates_internal, dim=-1)
        # ...
        combined_scores = topk_scores_1.unsqueeze(-1) + topk_scores_2.unsqueeze(-2)
        # ...
        combined_indices = (
            topk_indices_1.unsqueeze(-1) * self.num_keys + topk_indices_2.unsqueeze(-2)
        )
        # ...
        combined_scores = combined_scores.view(bsz, seq_len, -1)
        combined_indices = combined_indices.view(bsz, seq_len, -1)
        # ...
        candidate_scores, candidate_pk_indices = combined_scores.topk(
            self.num_candidates, dim=-1
        )
        candidate_indices = combined_indices.gather(-1, candidate_pk_indices)
        # ...
        candidate_scores = F.softmax(candidate_scores, dim=-1)
        candidate_scores = self.dropout(candidate_scores)

        return candidate_indices, candidate_scores
\end{lstlisting}

\begin{implnote}[title={In the code: where the sub-keys come from}]
\repofile{repos/SCUT-HCC__ExplicitLM/convert_conversations_to_labeled.py} embeds every knowledge
sentence with a sentence encoder (\code{BAAI/bge-base-en-v1.5}, normalized), then performs a
two-level \emph{residual quantization}: MiniBatch $k$-means with $\sqrt{N}$ clusters gives the row
sub-keys; $k$-means on the residuals (embedding minus its row centroid) gives the column
sub-keys. Entry $i$ is mapped to the grid cell (row cluster, residual cluster), and for each training
query the 32 nearest knowledge sentences (by exact inner-product search with FAISS) are stored as
targets. \code{knowledge_dim} $= 1536$ is twice the encoder's width, one half per sub-key set.
\code{train_router.py} then trains the query network with \code{compute_loss_soft}: the product-key
log-probability of a target, $\log p_1(u) + \log p_2(v)$ for grid cell $(u,v)$, weighted by a
softmax of the target similarities at temperature $0.1$ --- a soft-label cross-entropy.
\end{implnote}

\begin{implnote}[title={In the code: stage 2, fusion and gradients}]
In \code{Qwen3ExplicitLMBlock.forward} the query is the masked \emph{mean} of the normalized hidden
states over the sequence, so there is one candidate set per sequence and layer, broadcast to all
positions. The candidates' token ids are gathered from the bank and embedded with the tied
embedding table in chunks of 32, averaged over the $\ell_m$ tokens, and compared with each
position's normalized hidden state by cosine similarity. \code{gumbel_softmax_selection}
(\cref{lst:c10-explicit-gumbel}) implements \cref{eq:c10-gumbel} with $\tau_g$ =
\code{gumbel_temperature} $= 1.0$. \code{GatedMemoryFusion} implements \cref{eq:c10-explicit-fuse}.
Two consequences are worth knowing. First, the sequence mean contains future tokens, so the same
causality caveat as for LM2 applies in teacher-forced training, while generation (\code{_stream})
recomputes the whole prefix at each step and sees only the past. Second, after
\code{generate_candidates} the block uses only the candidate \emph{indices}; the candidate scores
are expanded but not used, so in \code{train_memory.py} the language-modeling loss does not reach
the query network through this path, even though its parameters are marked trainable. The coarse
stage is trained by the separate router objective above.
\end{implnote}

\begin{lstlisting}[caption={Excerpt from \texttt{repos/SCUT-HCC\_\_ExplicitLM/models/core/Qwen3ExplicitLMBlock.py} (body of \texttt{gumbel\_softmax\_selection}).},label={lst:c10-explicit-gumbel}]
        gumbel_noise = -torch.log(-torch.log(torch.rand_like(similarity_scores) + 1e-20) + 1e-20)
        logits = (similarity_scores + gumbel_noise) / temperature
        soft_weights = F.softmax(logits, dim=-1)
        if hard:
            _, max_indices = soft_weights.max(dim=-1, keepdim=True)
            hard_weights = torch.zeros_like(soft_weights).scatter_(-1, max_indices, 1.0)
            selection_weights = hard_weights - soft_weights.detach() + soft_weights
            selected_indices = max_indices.squeeze(-1)
        else:
            selection_weights = soft_weights
            selected_indices = torch.argmax(soft_weights, dim=-1)
        return selection_weights, selected_indices
\end{lstlisting}

\begin{implnote}[title={In the code: updating the bank at inference}]
\code{MemoryBankUpdater} (\repofile{repos/SCUT-HCC__ExplicitLM/utils/memory_bank_updater.py})
takes facts from \code{FactExtractor} (\repofile{repos/SCUT-HCC__ExplicitLM/utils/fact_extractor.py},
which wraps an LLMLingua-2 prompt compressor with compression rate $0.4$), splits the compressed
text at sentence boundaries, drops fragments of 10 characters or fewer, tokenizes and pads each
fact to $\ell_m$ tokens, and overwrites slots chosen by the configured strategy (\code{"fifo"} by
default; \code{"similarity"} currently falls back to random). The model also registers a
\code{freeze_mask} buffer (all \code{False}), and
\repofile{repos/SCUT-HCC__ExplicitLM/utils/dataset_processor/dataset_processor.py} keeps a
\code{FREEZE_RATIO = 0.2} that restricts evaluation questions to the first 20\% of database
entries --- remnants consistent with the paper's frozen/learnable partition. Note that overwriting
a slot's token ids changes its vector $\bar{\ve}_i$ but not the sub-keys of its grid cell, which
were derived from the \emph{old} content; a newly written fact is therefore reached through the
address of the slot it replaced, unless the keys are rebuilt.
\end{implnote}

\begin{lstlisting}[caption={Reference implementation (ours): two-stage retrieval with product keys and a straight-through Gumbel-Softmax (per-token queries).},label={lst:c10-explicit-ref}]
import torch, torch.nn.functional as F

def two_stage_retrieve(h, f_q, C1, C2, bank_ids, emb, k_short=128, c=16, tau_g=1.0):
    """h: (B, T, d) normalized hidden states; f_q: query network d -> dk;
    C1, C2: (n, dk/2) sub-keys; bank_ids: (n*n, l_m) token ids; emb: nn.Embedding."""
    n = C1.shape[0]
    q = f_q(h)                                               # (B, T, dk), one query per token
    q1, q2 = q.chunk(2, dim=-1)
    a = F.normalize(q1, dim=-1) @ F.normalize(C1, dim=-1).T  # (B, T, n) cosine scores
    b = F.normalize(q2, dim=-1) @ F.normalize(C2, dim=-1).T
    av, ai = a.topk(k_short, -1)
    bv, bi = b.topk(k_short, -1)
    grid = (av[..., :, None] + bv[..., None, :]).flatten(-2)
    gid = (ai[..., :, None] * n + bi[..., None, :]).flatten(-2)
    cand = gid.gather(-1, grid.topk(c, -1).indices)          # (B, T, c) entry ids
    e_bar = emb(bank_ids[cand]).mean(-2)                     # (B, T, c, d): mean token embedding
    u = F.cosine_similarity(h[..., None, :], e_bar, dim=-1)  # (B, T, c) fine scores
    g = -torch.log(-torch.log(torch.rand_like(u).clamp_min(1e-20)))
    soft = torch.softmax((u + g) / tau_g, dim=-1)
    hard = F.one_hot(soft.argmax(-1), c).to(soft.dtype)
    y = hard - soft.detach() + soft                          # straight-through weights
    m = (y[..., None] * e_bar).sum(-2)                       # (B, T, d) selected memory
    sim_loss = -(y * u).sum(-1).mean()                       # auxiliary similarity loss
    return m, cand.gather(-1, hard.argmax(-1, keepdim=True)), sim_loss
\end{lstlisting}

\begin{results}
\textbf{Source files:} \repofile{repos/SCUT-HCC__ExplicitLM/README.md}; survey PDF. The repository
reports no evaluation numbers, and the survey describes ExplicitLM qualitatively. The paper's
abstract reports substantial improvements on knowledge-intensive tasks over a standard Transformer,
larger gains in low-data regimes, and a strong correlation between memory retrieval and
performance (correct predictions retrieve relevant entries more often). See the paper for the
numbers.
\end{results}

\subsection{Discussion}

ExplicitLM's distinctive property is that the memory's content is text. A reader can decode which
entry a layer selected for a token, and an operator can correct a fact by writing new token ids ---
an operation that is hard to perform on entangled parametric memory and that the survey singles
out as the significance of editable memory banks. The product-key stage makes a million entries
affordable, and the straight-through Gumbel-Softmax keeps the selection discrete without giving up
gradients for the fine stage.

The price is representational. An entry's vector is the mean of its token embeddings, which
discards word order (``A founded B'' and ``B founded A'' coincide) and limits how much a 16-token
entry can say; one entry is selected per token and layer; and the address structure (sub-keys built
by clustering sentence embeddings) is separate from the content, so edits and keys can drift
apart. In the repository version the backbone is frozen and the coarse router is trained against a
sentence-encoder teacher, which turns the system into a learned retriever with fusion rather than
a jointly trained memory; whether the paper's version differs in these respects requires the paper.
Compared with Engram (\cref{sec:c10-engram}), ExplicitLM trades a static, collision-prone but
$\bigO(1)$ address for a content-based, readable but more expensive one; compared with kNN-LM
(\cref{sec:c10-knnlm}), its memory is consumed inside every layer rather than mixed into the output
distribution.

% =====================================================================================
\section{MemoryLLM: feed-forward layers as plug-and-play token memory}\label{sec:c10-memoryllm}

\paragraph{Which MemoryLLM?} Two different systems share the name. \emph{MEMORYLLM}
\citep{wang2024memoryllm} (2024) is the self-updatable model with a latent memory pool described in
\cref{sec:c10-precursors}; it is online and its code is public. \emph{MemoryLLM: Plug-n-Play
Interpretable Feed-Forward Memory for Transformers} \citep{jaiswal2026memoryllm} (Apple, ICML 2026)
is the offline design of this section. No public code for it was found, so we explain it from the
survey's description and the authors' public summary, and we mark our own formalization.

\subsection{Motivation}

\Cref{eq:c10-ffn-kv} says that an FFN is a key--value memory, and FFNs hold most of a Transformer's
parameters: with attention projections $\mW_Q,\mW_K,\mW_V,\mW_O \in \R^{d\times d}$ ($4d^2$
parameters) and an FFN with $d_{\mathrm{ff}} = 4d$ and two matrices ($8d^2$) --- or a SwiGLU FFN
with $d_{\mathrm{ff}} = \tfrac{8}{3}d$ and three matrices, again $8d^2$ --- the FFN is two thirds of
each block. Yet this memory is hard to study, because its input is the residual stream: a mixture
of the token, its context and everything earlier layers computed. One cannot say which ``memory
locations'' a given token reads, because the answer depends on the whole context.

MemoryLLM removes that dependence. According to the survey (\S IV-B(e)), it ``trains decoupled FFN
memories directly from token embeddings, enabling their outputs to be precomputed as token-level
lookups''. The FFNs become context-free, token-wise retrieval modules; the attention layers keep
all the context processing. Two benefits follow: interpretability (the memory accessed by each
vocabulary item can be enumerated) and efficiency (FFN outputs can be precomputed, stored outside
GPU memory and fetched on demand). A ``flexible'' variant sits between this design and a
conventional Transformer.

\begin{taxonomy}
\textbf{Representation:} explicit --- the survey places it under lookup-based memory, as feed-forward
layers reinterpreted as context-free token-wise retrieval modules \citep[\S IV-B(e)]{zhoubian2026memory};
Table~I's paradigm label is ``Context-Free Feed-Forward Lookup''.
\textbf{Update dynamics:} offline. \textbf{Persistence:} long-term.
\textbf{Update rule:} not listed in Table~II; in our reading objective-induced/structural (the
memory is fixed by training and by an offline precomputation).
\end{taxonomy}

\subsection{Formulation}

\paragraph{Decoupling the FFN from the residual stream.} A standard pre-norm block computes
\begin{equation}
  \mX' = \mX^{(\ell)} + \mathrm{Attn}^{(\ell)}\bigl(\RMSNorm(\mX^{(\ell)})\bigr),\qquad
  \mX^{(\ell+1)} = \mX' + \mathrm{FFN}^{(\ell)}\bigl(\RMSNorm(\mX')\bigr).
\end{equation}
Public summaries of MemoryLLM describe the block as
$\mX^{(\ell+1)} = \mX^{(\ell)} + \mathrm{Attn}(\mX^{(\ell)}) + \mathrm{FFN}(\mX^{(0)})$, where
$\mX^{(0)}$ holds the token embeddings $\mathrm{Emb}(x_t)$. One way to formalize this, with the
normalizations written out, is
\begin{equation}
  \mX^{(\ell+1)} = \mX^{(\ell)} + \mathrm{Attn}^{(\ell)}\bigl(\RMSNorm(\mX^{(\ell)})\bigr)
  + \mathrm{FFN}^{(\ell)}\bigl(\RMSNorm(\mX^{(0)})\bigr),
  \label{eq:c10-memoryllm}
\end{equation}
so attention and FFN act in parallel on different inputs: attention on the contextual stream, the
FFN on the context-free token embedding. The paper's exact parameterization (placement of norms,
embedding scaling, whether FFN inputs are shared across layers) may differ.

\paragraph{Token-wise lookups.} Both $\RMSNorm$ and the FFN act on each row separately, so row $t$
of $\mathrm{FFN}^{(\ell)}(\RMSNorm(\mX^{(0)}))$ is $\mathrm{FFN}^{(\ell)}(\RMSNorm(\mathrm{Emb}(x_t)))$,
a function of the token id $x_t$ and the layer index alone. Define the table of token-wise lookups
(ToLs)
\begin{equation}
  \mathcal{T}[v, \ell] = \mathrm{FFN}^{(\ell)}\bigl(\RMSNorm(\mathrm{Emb}(v))\bigr) \in \R^{d},
  \qquad v \in \mathcal{V},\; \ell = 1,\dots,L .
  \label{eq:c10-tol}
\end{equation}

\begin{proposition}[ToLs are exact]\label{prop:c10-tol}
For every input sequence, replacing the FFN term of \cref{eq:c10-memoryllm} by the row
$\mathcal{T}[x_t, \ell]$ at every position $t$ and layer $\ell$ leaves all hidden states and logits
unchanged.
\end{proposition}
\begin{proof}
At every layer the two computations add the same vector to row $t$, because
\[
  \bigl[\mathrm{FFN}^{(\ell)}(\RMSNorm(\mX^{(0)}))\bigr]_t
  = \mathrm{FFN}^{(\ell)}(\RMSNorm(\mathrm{Emb}(x_t))) = \mathcal{T}[x_t,\ell];
\]
everything else in the recursion is identical. Induction over $\ell$ gives equality of all hidden states.
\end{proof}

The proposition fails for a conventional block, whose FFN input $\RMSNorm(\mX')$ depends on the
context; this is the whole point of the decoupling. Our NumPy check (scratch space) built a small
random model of the form \eqref{eq:c10-memoryllm}, precomputed the table, and obtained logits
identical to the online computation (maximum difference exactly 0).

\paragraph{The price in storage.} The table has $\abs{\mathcal{V}}\, L\, d$ entries, against
$L \cdot 3 d\, d_{\mathrm{ff}}$ parameters for SwiGLU FFNs, a ratio of
$\abs{\mathcal{V}}/(3 d_{\mathrm{ff}})$. For illustrative values $\abs{\mathcal{V}} = 131{,}072$,
$d = 4096$, $d_{\mathrm{ff}} = 8192$ and $L = 32$ (our numbers, not the paper's), the table holds
$1.7 \times 10^{10}$ numbers --- about five times the FFN parameters --- and each token needs a
$L d$-vector, $256$\,KB in 16-bit precision. ToLs therefore trade FLOPs for capacity in slower
memory: they are worth it when the table lives off the GPU (host memory or flash) and is fetched by
token id, which, as for Engram, is known before the forward pass and highly skewed (frequent
tokens can be cached). The abstract describes exactly this on-demand transfer between VRAM and
storage.

\paragraph{Interpretability.} With a context-free input, key $i$ of layer $\ell$ responds to token
$v$ with activation $\phi(\vk_i\T\RMSNorm(\mathrm{Emb}(v)))$, a finite list that can be computed
for the whole vocabulary. The set of memory locations a token reads is fixed, so one can ask
directly which tokens share keys and what each value writes --- the analysis the paper's title
promises. In a conventional Transformer the same question has a different answer in every context.

\paragraph{Flex-MemoryLLM.} Removing context from every FFN costs quality, and the authors propose
Flex-MemoryLLM to bridge the gap. Public summaries describe a split of the FFN parameter budget into
a small context-aware module FFN-C that acts on the residual stream and a larger context-free
module FFN-M that acts on the token embeddings and can be offloaded. One way to formalize this is
\begin{equation}
\begin{aligned}
  \mX^{(\ell+1)} = \mX^{(\ell)} &+ \mathrm{Attn}^{(\ell)}\bigl(\RMSNorm(\mX^{(\ell)})\bigr)
  + \mathrm{FFN}^{(\ell)}_{\mathrm{C}}\bigl(\RMSNorm(\mX^{(\ell)})\bigr)\\
  &+ \mathrm{FFN}^{(\ell)}_{\mathrm{M}}\bigl(\RMSNorm(\mX^{(0)})\bigr),
\end{aligned}
  \label{eq:c10-flex}
\end{equation}
with a fraction $\beta$ of the FFN parameters in FFN-C. $\beta = 0$ recovers
\cref{eq:c10-memoryllm} and $\beta = 1$ a conventional parallel block; only FFN-M can be tabulated,
and \cref{prop:c10-tol} applies to it unchanged (our NumPy check covered this case as well). The
paper's exact parameterization may differ.

\subsection{Algorithms and complexity}

\begin{algorithm}[t]
\caption{Building and using token-wise lookups (our formalization)}\label{alg:c10-tol}
\begin{algorithmic}[1]
\Require trained model of the form \eqref{eq:c10-memoryllm} or \eqref{eq:c10-flex}
\For{$v \in \mathcal{V}$ and $\ell = 1,\dots,L$} \Comment{offline, once}
  \State $\mathcal{T}[v,\ell] \gets \mathrm{FFN}^{(\ell)}_{(\mathrm{M})}(\RMSNorm(\mathrm{Emb}(v)))$
\EndFor
\State store $\mathcal{T}$ in host memory or on flash; optionally cache rows of frequent tokens on the GPU
\Statex \textit{Inference on tokens $x_{1:T}$:}
\State prefetch rows $\mathcal{T}[x_t, 1{:}L]$ as soon as $x_t$ is known
\For{$\ell = 1,\dots,L$}
  \State $\mX \gets \mX + \mathrm{Attn}^{(\ell)}(\RMSNorm(\mX)) \;[+\, \mathrm{FFN}^{(\ell)}_{\mathrm{C}}(\RMSNorm(\mX))] + \mathcal{T}[x_{1:T}, \ell]$
\EndFor
\end{algorithmic}
\end{algorithm}

Precomputation costs $\abs{\mathcal{V}} L$ FFN evaluations, once. At inference the FFN-M FLOPs
vanish and are replaced by $L d$ numbers of memory traffic per token; attention (and FFN-C) are
unchanged, so the KV cache and the context-dependent cost are those of the backbone.

\begin{lstlisting}[caption={Reference implementation (ours): a MemoryLLM-style block and the token-wise lookup table.},label={lst:c10-memllm-ref}]
import torch, torch.nn as nn, torch.nn.functional as F

class SwiGLU(nn.Module):
    def __init__(self, d, d_ff):
        super().__init__()
        self.w1, self.w3 = nn.Linear(d, d_ff, bias=False), nn.Linear(d, d_ff, bias=False)
        self.w2 = nn.Linear(d_ff, d, bias=False)
    def forward(self, x):
        return self.w2(F.silu(self.w1(x)) * self.w3(x))

class MemoryBlock(nn.Module):
    """x_next = x + Attn(norm(x)) [+ FFN_C(norm(x))] + FFN_M(norm(x0))."""
    def __init__(self, d, n_heads, d_ff_mem, d_ff_ctx=0):
        super().__init__()
        self.n_a, self.n_m = nn.RMSNorm(d), nn.RMSNorm(d)
        self.attn = nn.MultiheadAttention(d, n_heads, bias=False, batch_first=True)
        self.ffn_m = SwiGLU(d, d_ff_mem)                       # context-free memory
        self.ffn_c = SwiGLU(d, d_ff_ctx) if d_ff_ctx else None  # Flex variant
    def forward(self, x, x0=None, tol_rows=None):
        T = x.shape[1]
        causal = torch.ones(T, T, dtype=torch.bool, device=x.device).triu(1)
        z = self.n_a(x)
        out = x + self.attn(z, z, z, attn_mask=causal, need_weights=False)[0]
        if self.ffn_c is not None:
            out = out + self.ffn_c(z)
        mem = tol_rows if tol_rows is not None else self.ffn_m(self.n_m(x0))
        return out + mem

@torch.no_grad()
def build_tol(blocks, emb):                  # emb: nn.Embedding(|V|, d)
    e = emb.weight                           # (|V|, d), every token once
    return torch.stack([b.ffn_m(b.n_m(e)) for b in blocks], dim=1)   # (|V|, L, d)

def forward_with_tol(blocks, emb, tol, ids):
    x = emb(ids)
    for l, b in enumerate(blocks):
        x = b(x, tol_rows=tol[ids, l])       # a gather replaces the FFN-M compute
    return x
\end{lstlisting}

\begin{results}
\textbf{Source files:} survey PDF (no repository exists for this paper). The survey reports no
numbers; it describes MemoryLLM as weakening the dependence of FFN memory access on the attention
context and turning part of the Transformer into an explicit storage-and-lookup substrate
(\S IV-B(e)). The paper's abstract states that the context-free FFNs can be precomputed as
token-wise lookups that are transferred on demand between VRAM and storage, improving inference
efficiency, and that Flex-MemoryLLM narrows the performance gap caused by training FFNs on
context-free token embeddings. See the paper for the numbers.
\end{results}

\subsection{Discussion}

MemoryLLM is the most radical of the four designs: it does not add a memory to the Transformer,
it reinterprets an existing component. The survey calls it a bridge between parameter memory and
retrieval memory --- the knowledge is still in neural parameters, but its access pattern is that of
a structured lookup table (\S IV-B(e)). Its relation to Engram is instructive. Both make part of
the network token-indexed, deterministic and offloadable. Engram keys on short $n$-grams, inserts
memory at a few layers, adds it to a full Transformer and gates it by the hidden state; MemoryLLM
keys on single tokens, puts a memory in every layer, and replaces the contextual FFN by it, with
no gate in the base design. A MemoryLLM layer is in this sense a ``unigram Engram without gating''
--- a comparison that is ours, not the authors'.

The cost of the design is expressiveness: in \cref{eq:c10-memoryllm} the only context-dependent
nonlinearities are those of attention and normalization, so contextual feature computation that a
conventional FFN performs must move into attention or into FFN-C. Flex-MemoryLLM exposes this
trade-off as a single knob. Updates remain offline: changing what a token ``knows'' means editing
its ToL rows, which is easy mechanically (the rows are addressed by token id) but has no learning
rule attached.

% =====================================================================================
\section{Comparison and discussion}\label{sec:c10-comparison}

\Cref{tab:c10-comparison} places the four systems next to the three reference designs of the
background section. Read it along three questions, followed by a lesson about the code.

\begin{table}[t]
\centering
\footnotesize
\caption{Slot and lookup memories compared. $N$: entries; $d$: model width; $k$: entries read; $T$:
sequence length; $L$: layers. ``Read cost'' is per token and per memory layer.}
\label{tab:c10-comparison}
\setlength{\tabcolsep}{3pt}
\begin{tabularx}{\linewidth}{@{}>{\raggedright\arraybackslash}p{1.85cm}
  >{\raggedright\arraybackslash}X >{\raggedright\arraybackslash}X
  >{\raggedright\arraybackslash}X >{\raggedright\arraybackslash}X
  >{\raggedright\arraybackslash}p{2.1cm}@{}}
\toprule
System & What is stored & Addressing & Capacity scales with & Update mode & Read cost \\
\midrule
kNN-LM & (context vector, next token) pairs from a corpus & nearest neighbours of the final hidden state & corpus size & offline, append/edit without training & ANN search, sublinear in $N$ \\
PKM & $N$ value vectors, factorized keys & product-key top-$k$ of a query from $\vh_t$ & $N = n^2$ values & offline (gradient, sparse) & $\bigO(\sqrt{N} d_k + k^2 + k d)$ \\
MoE & $E$ expert MLPs & learned router on $\vh_t$ & number of experts & offline (gradient) & $k$ expert MLPs \\
\midrule
LM2 & $N$ slot vectors per sequence (bank) & cross-attention from tokens & $N d$, tied to $N = T$ in code & gated rewrite every forward pass; offline in survey & $\bigO(N d)$ ($\bigO(TNd)$ per layer) \\
Engram & hashed $n$-gram embeddings & hash of last $\le n_{\max}$ token ids; scalar gate on $\vh_t$ & table rows (billions of parameters), no FLOP increase & offline (gradient; tables sharded) & $(n_{\max}{-}1)K$ row reads $+\,\bigO(R d\, d_{\mathrm{mem}})$ \\
ExplicitLM & $N$ token sequences (readable ids) & product keys (coarse) + cosine/Gumbel (fine) & $N$ entries; values share the embedding table & offline; repo rewrites entries at inference & $\bigO(\sqrt{N}d_k + k'^2 + c\,\ell_m d)$ \\
MemoryLLM & FFN outputs per token and layer (ToLs) & token id & $\abs{\mathcal{V}} L d$ table & offline (training + precomputation) & $L d$ memory reads, no FFN FLOPs \\
\bottomrule
\end{tabularx}
\end{table}

\paragraph{Content addressing versus id addressing.} LM2, ExplicitLM, PKM and MoE address memory
with a vector computed from the hidden state. That makes the address context-sensitive but serial:
nothing can be fetched before the layer runs. Engram and MemoryLLM address memory with token ids.
That makes the address blind to context but \emph{static}, so it can be computed ahead of time,
prefetched across PCIe, sharded, cached by frequency and moved down the memory hierarchy. Engram
reintroduces context after the lookup with its gate; MemoryLLM's base design does not, and its
flexible variant reintroduces context through a separate FFN. Hybrid designs that route between
memories adaptively are the subject of \cref{ch:routing}.

\paragraph{Where capacity lives and what it costs.} LM2's memory is per-sequence working memory
whose size is tied to the sequence and whose read is as expensive as attention. The three lookup
memories put capacity into persistent tables whose size is decoupled from compute: hundreds of
millions to billions of rows for Engram, a million readable entries for ExplicitLM (with no value
parameters at all), $\abs{\mathcal{V}} L$ rows for MemoryLLM. Engram's allocation study
quantifies how such capacity should be traded against MoE capacity under a fixed budget.

\paragraph{Writing and editing.} Only LM2 writes during the forward pass, and its writes are soft,
dense and opaque. ExplicitLM is the only design whose entries can be read and rewritten by a human;
its repository even implements inference-time rewriting, at the price of keys that do not follow
the new content. Engram's and MemoryLLM's memories are fixed after training; they are long-term
stores, not working memories. None of the four uses the optimization-based writes of the
test-time-training family (\cref{ch:ttt}), and none addresses the stability--plasticity questions
that the multi-timescale methods of \cref{ch:multiscale} take up.

\paragraph{Implementation lessons.} Two of the official code bases pool a statistic over the whole
sequence (LM2's input gate, ExplicitLM's retrieval query), which breaks causality under teacher
forcing; and the public Engram code is a demonstration that differs from the paper in small but
real ways (signed square root in the gate, no zero-initialized convolution). Reading the code
closely is part of understanding these systems. \Cref{ch:synthesis} returns to how slot and lookup
memories combine with the recurrent and attention memories of the earlier chapters.

\begin{summarybox}
\begin{itemize}
  \item Explicit, addressable memory stores knowledge in entries that persist, are individually
        addressed, and scale in number independently of the hidden size. The address step decides
        the cost: dense attention $\bigO(Nd)$, product keys $\bigO(\sqrt{N}d + k^2)$ (exact by
        \cref{prop:c10-pkm}), hashing and token ids $\bigO(1)$.
  \item FFNs are key--value memories with dense, context-dependent addresses; MoE sparsifies the
        address of whole functions (conditional computation); hashed tables sparsify the address of
        stored vectors (conditional memory).
  \item \textbf{LM2} adds to every decoder layer a bank of slots, read by cross-attention and
        rewritten by LSTM-style input/forget gates, with gated output into the residual stream; in
        the code the bank is threaded through the layers and carried across batches.
  \item \textbf{Engram} retrieves hashed $2$- and $3$-gram embeddings in $\bigO(1)$, gates them with
        the hidden state and fuses them through a short causal convolution. Under fixed parameters
        and FLOPs, splitting the sparse budget between MoE and Engram beats pure MoE (U-shaped
        allocation curve), and deterministic addresses allow tables in host memory with small
        throughput loss.
  \item \textbf{ExplicitLM} keeps a million human-readable token sequences, finds candidates with
        product keys and selects one with a straight-through Gumbel-Softmax; its entries can be
        inspected and rewritten.
  \item \textbf{MemoryLLM} feeds FFNs with token embeddings only, so their outputs become
        token-wise lookup tables that can be precomputed and offloaded; Flex-MemoryLLM keeps a small
        contextual FFN.
  \item Pitfall: any gate or query pooled over the whole sequence breaks causality in
        autoregressive training.
\end{itemize}
\end{summarybox}

\section*{Exercises}

\begin{exercise}[Product keys with short lists]\label{exr:c10-pkm}
(a) Extend \cref{prop:c10-pkm} to scores with ties: show that \emph{some} global top-$k$ set lies
in $\mathcal{I}_1 \times \mathcal{I}_2$ when ties are broken consistently. (b) Find the smallest
example ($n = 2$, $k = 2$, shortlist $k' = 1$) in which the shortlist misses a global top-$k$ pair,
for instance $\va = (1, 0)$, $\vb = (1, 0.9)$. (c) Write the cost of a two-level ``product of
three'' key (three sub-key sets of size $N^{1/3}$) and explain why the candidate merge becomes the
bottleneck.
\end{exercise}

\begin{exercise}[Collisions under multi-head hashing]\label{exr:c10-hash}
Assume a random-function model for $K$ independent heads, each with $P$ rows, and $U$ distinct
$n$-grams. (a) Show that the expected fraction of $n$-grams sharing their row in a given head is
$1 - (1-1/P)^{U-1}$, and evaluate it for $U/P = 4$. (b) Show that the probability that a fixed
$n$-gram has the \emph{same partner} in all $K$ heads is at most $(U-1)/P^K$. (c) Now use
\cref{prop:c10-crt} to explain why, for Engram's deterministic hash with distinct primes whose
product exceeds $2^{63}$, the answer to (b) is governed by collisions of the 63-bit mixed value and
not by $P^K$.
\end{exercise}

\begin{exercise}[A causal LM2]\label{exr:c10-lm2}
(a) Take $T = N = d = 1$ and write the code's gate (mean over positions of $\ReLU(w e_t)$, here of a
two-token sequence) explicitly; show that the output at position 1 depends on the token at
position 2. (b) Modify \cref{lst:c10-lm2-ref} so that the bank written during segment $s$ is read
only by segment $s+1$ (segment-level recurrence) and prove the result is causal for any $N$ (no
longer tied to $T$). (c) Give the time per generated token and the state size of your design in
decoding.
\end{exercise}

\begin{exercise}[Gumbel maxima]\label{exr:c10-gumbel}
(a) Show that $\max_i (u_i + G_i)$, with $G_i$ i.i.d.\ standard Gumbel, is Gumbel-distributed with
location $\log\sum_i e^{u_i}$. (b) Deduce $\E[\max_i(u_i + G_i)] = \log\sum_i e^{u_i} + \gamma_{\mathrm{E}}$,
where $\gamma_{\mathrm{E}} \approx 0.5772$ is the Euler--Mascheroni constant (the mean of a standard
Gumbel). (c) Check both statements by simulation in NumPy, and explain why
\cref{lst:c10-explicit-ref} clamps the uniform sample before taking logarithms.
\end{exercise}

\begin{exercise}[Engram's sparsity budget]\label{exr:c10-alloc}
Use the numbers of \repofile{repos/deepseek-ai__Engram/Engram_paper.pdf} (Section~4.1 and Table~1):
MoE-27B has 72 routed experts (top-6) and 2 shared experts; Engram-27B has 55 routed experts and a
5.7B-parameter Engram at $\rho = 74.3\%$; both have the same $P_{\mathrm{tot}}$ and
$P_{\mathrm{act}}$. (a) From \cref{eq:c10-rho} compute $P_{\mathrm{sparse}}$ and the parameter
count of one routed expert summed over all MoE layers. (b) Check that removing 17 experts frees about
5.7B parameters, and that the inactive experts of MoE-27B account for $P_{\mathrm{sparse}}$.
(c) Estimate the per-token FLOPs of the two Engram layers from the configuration in Appendix~A and
compare them with $2 P_{\mathrm{act}}$.
\end{exercise}

\begin{exercise}[Token-wise lookups in practice]\label{exr:c10-tol}
For a MemoryLLM-style model with $\abs{\mathcal{V}} = 151{,}936$, $L = 36$, $d = 2560$ and SwiGLU
FFNs with $d_{\mathrm{ff}} = 9728$: (a) compute the ToL size in bytes at 16-bit precision and compare
it with the FFN parameter count; (b) compute the bytes fetched per token and the PCIe bandwidth
needed at 50 tokens per second without caching; (c) assuming token frequencies follow a Zipf law
with exponent 1, estimate what fraction of fetches a GPU cache of the 10{,}000 most frequent tokens
absorbs; (d) implement \cref{lst:c10-memllm-ref} and verify \cref{prop:c10-tol} numerically,
including the Flex variant.
\end{exercise}
